\chapter{Baseband Transmission and Pulse Shaping}
\label{ch:baseband}

\begin{chapterintro}
Every digital link, from a 1960s T1 span to a 224\,Gb/s lane inside an AI data centre, must
turn a sequence of bits into a physical waveform, push it through a channel that blurs and
adds noise, and turn it back into bits. This chapter covers that problem in its simplest
setting, \emph{baseband}, where no carrier is involved. It starts with line codes (the
practical rules that make a waveform DC-free, self-clocking and error-checkable) and the
power spectra they produce. It then derives the Nyquist criterion for zero intersymbol
interference and the raised-cosine pulses that meet it, proves that the matched filter is
the best possible receiver filter, and works out the error probability of binary signalling
in Gaussian noise. The rest of the chapter covers the tools of the high-speed-link engineer:
eye diagrams, jitter and bathtub curves, channel loss and its equalisation, multilevel PAM
in today's 400G and 800G Ethernet, PCIe~6 and graphics memory, and partial-response
signalling from duobinary to the PRML read channels that made the hard-disk revolution
possible. It closes with a look at faster-than-Nyquist signalling. Everything in later
chapters on carrier modulation rests on the ideas developed here.
\end{chapterintro}

%======================================================================
\section{From bits to waveforms}
\label{sec:ch08:pam}
%======================================================================

A wire, a fibre and a radio channel share one property: none of them can carry a bit. They
carry voltages, currents, photons and fields, all continuous functions of time. The first
job of any digital transmitter is to choose a waveform for each group of bits, and the
first job of the receiver is to decide, from a noisy and distorted version of that waveform,
which group was sent. This chapter covers the simplest and most important family of
choices.

\subsection{The PAM model}

In \term{pulse-amplitude modulation} (PAM) the transmitter sends one pulse shape $p(t)$ every
$T$ seconds and scales each pulse by a symbol $a_k$ drawn from a small alphabet:
\begin{equation}
s(t)=\sum_{k=-\infty}^{\infty}a_k\,p(t-kT).
\label{eq:ch08:pam}
\end{equation}
With the binary alphabet $a_k\in\{-1,+1\}$ and a rectangular pulse, this is the familiar
square wave of polar signalling. With the alphabet $\{-3,-1,+1,+3\}$ it is 4-PAM. With a
complex alphabet (QPSK, 16-QAM) and a carrier, it becomes the quadrature modulation of
Chapter~\ref{ch:modulation}. Nearly every single-carrier digital system in this book,
including the downlink of every geostationary TV satellite, the upstream of a cable modem
and the electrical lanes inside an Ethernet switch, is an instance of~\eqref{eq:ch08:pam}.
OFDM (Chapter~\ref{ch:ofdm}) is a set of many PAM-like signals running in parallel.

Equation~\eqref{eq:ch08:pam} splits the design into two nearly independent problems.
Choosing the \emph{alphabet} and the mapping from bits to symbols sets the spectral
efficiency and the noise tolerance. Choosing the \emph{pulse} $p(t)$ sets the bandwidth,
the sensitivity to timing errors and whether neighbouring symbols interfere. Most of this
chapter is about the pulse; Chapter~\ref{ch:modulation} is about the alphabet.

\begin{figure}[htbp]\centering
\begin{tikzpicture}[node distance=0.42cm and 0.42cm,
  blk/.style={draw=navy,fill=navy!6,thick,minimum height=0.95cm,minimum width=1.35cm,align=center,font=\sffamily\scriptsize},
  arr/.style={-{Stealth[length=2mm]},thick,navy}]
\node[blk] (map) {Bits to\\symbols\\$a_k$};
\node[blk,right=of map] (tx) {Transmit\\filter\\$g_T(t)$};
\node[blk,right=of tx,fill=accent!6,draw=accent] (ch) {Channel\\$c(t)$};
\node[draw=accent,circle,thick,right=of ch,minimum size=0.6cm,font=\sffamily\small] (sum) {$+$};
\node[blk,right=of sum] (rx) {Receive\\filter\\$g_R(t)$};
\node[blk,right=of rx] (smp) {Sample\\at\\$t_0+kT$};
\node[blk,right=of smp] (dec) {Decide\\$\hat a_k$};
\node[above=0.45cm of sum,font=\sffamily\scriptsize,text=accent] (n) {noise $n(t)$};
\draw[arr] (map) -- (tx); \draw[arr] (tx) -- node[above,font=\scriptsize]{$s(t)$} (ch);
\draw[arr] (ch) -- (sum); \draw[arr] (sum) -- node[above,font=\scriptsize]{$r(t)$} (rx);
\draw[arr] (rx) -- node[above,font=\scriptsize]{$y(t)$} (smp); \draw[arr] (smp) -- node[above,font=\scriptsize]{$y_k$} (dec);
\draw[arr,accent] (n) -- (sum);
\draw[decorate,decoration={brace,amplitude=5pt,mirror},navy!70] ([yshift=-3pt]tx.south west) -- ([yshift=-3pt]rx.south east)
  node[midway,below=6pt,font=\sffamily\scriptsize,text=navy]{overall pulse $x(t)=g_T*c*g_R$: Nyquist criterion applies here};
\end{tikzpicture}
\caption{The baseband PAM link. The transmit filter, channel and receive filter combine
into a single end-to-end pulse $x(t)$; the receive filter also sets how much noise reaches
the sampler. The two design questions of this chapter are how to make $x(t)$ free of
intersymbol interference and how to split the filtering to maximise the signal-to-noise
ratio at the sampler.}
\label{fig:ch08_pam_chain}
\end{figure}

Figure~\ref{fig:ch08_pam_chain} shows the complete link. The transmitted pulse is the
response $g_T(t)$ of a transmit filter to an impulse of weight $a_k$; the channel adds its
own impulse response $c(t)$ and noise; the receiver filters with $g_R(t)$, samples once per
symbol at time $t_0+kT$ and decides. The end-to-end pulse seen at the sampler is
\begin{equation}
x(t)=g_T(t)*c(t)*g_R(t),
\label{eq:ch08:overall}
\end{equation}
and the sampler output is
\begin{equation}
y_k=a_k\,x(t_0)+\underbrace{\sum_{m\ne k}a_m\,x\big(t_0+(k-m)T\big)}_{\text{intersymbol interference}}+\;\nu_k,
\label{eq:ch08:sample}
\end{equation}
where $\nu_k$ is filtered noise. The first term is the wanted symbol, the last is noise, and
the middle term, \term{intersymbol interference} (ISI), is the signal getting in its own way.
Sections~\ref{sec:ch08:nyquist} and~\ref{sec:ch08:mf} show how to eliminate the middle term
and minimise the last.

\subsection{Symbol rate, bit rate and the baud}

If each symbol carries $b=\log_2 M$ bits from an $M$-ary alphabet, the \term{bit rate} is
$R_b=b/T$ and the \term{symbol rate} is $R_s=1/T$. Symbol rate is measured in
\term{baud} (Bd), after \'Emile Baudot, the French telegraph engineer whose five-unit
printing-telegraph code of the 1870s was the ancestor of every character code since. The
unit was adopted by the international telegraph community in the 1920s. One baud is one
symbol per second, and it is the symbol rate, not the bit rate, that sets the bandwidth a
signal needs. That is why a 100\,Gb/s Ethernet lane that uses 4-PAM runs at 53.125\,GBd,
and why engineers who say ``a 9600-baud modem'' about a V.32 modem that sent 2400 symbols
per second at 4 bits per symbol were, strictly, wrong by a factor of four.

\begin{keyidea}[Bandwidth follows the symbol rate; noise tolerance follows the alphabet]
A PAM signal needs a bandwidth of at least $R_s/2$\,Hz (Section~\ref{sec:ch08:nyquist}),
whatever the alphabet. Doubling the bits per symbol at a fixed symbol rate doubles the bit
rate in the same bandwidth but packs the levels closer together, so the receiver needs more
signal-to-noise ratio. Every trade-off in this chapter---NRZ versus PAM-4, binary versus
duobinary, rectangular versus raised-cosine---is a variation of this exchange between
bandwidth, power and complexity.
\end{keyidea}

\begin{history}[The telegraph pulse problem]
Intersymbol interference was the first great problem of electrical communication. The first
transatlantic cable of 1858 was a single conductor wrapped in gutta-percha and armoured in
iron, about 3000\,km long. Its distributed resistance and capacitance smeared each
telegraph pulse into a slow rise and a slower tail lasting many seconds, so that
consecutive dots and dashes ran together. William Thomson (later Lord Kelvin) had analysed
this in 1854-55 and showed that the smearing time grows with the \emph{square} of the cable
length, his ``law of squares''. He designed the sensitive mirror galvanometer that let the
1858 cable work at all, at a rate of roughly a word or less per minute. The cable failed
after a few weeks, its insulation damaged by the high voltages the chief electrician, Wildman
Whitehouse, applied in an attempt to push the signals through. The 1866 cable, built on
Thomson's principles, worked at around eight words per minute. Oliver Heaviside's
telegrapher's equations of the 1880s later explained that adding series inductance
(``loading'') could make a line distortionless, an idea Michael Pupin and George Campbell
turned into loading coils around 1900. The telegraph engineer's question---how fast can
pulses be sent before they merge?---is exactly the question this chapter answers.
\end{history}

%======================================================================
\section{Line codes}
\label{sec:ch08:linecodes}
%======================================================================

Before the theory of optimum pulses, consider the practical question every wireline engineer
faces first: which waveform to put on the wire for each bit. A \term{line code} is a rule
that maps the bit stream to a sequence of levels, usually with rectangular pulses. The
choice is driven not by the Gaussian-noise optimum but by the physical constraints of real
cables, transformers, fibres and clock-recovery circuits.

\subsection{What a line code must do}

A good line code for a practical link should provide most of the following.
\begin{description}[style=nextline,leftmargin=1.2em,font=\sffamily\bfseries\small\color{navy}]
\item[DC balance.] Most wireline links are AC-coupled through transformers (telephone lines,
Ethernet over twisted pair) or series capacitors (every high-speed SerDes), so the channel
cannot pass DC. A signal with a long-term average away from zero causes \term{baseline
wander}: the coupling capacitor charges, the decision threshold drifts, and errors follow.
The code should have no spectral content at, or near, zero frequency.
\item[Clock content.] The receiver must recover the transmitter's clock from the data
(Chapter~\ref{ch:sync}). Clock recovery circuits learn timing from transitions, so the code
must guarantee a minimum \term{transition density} and a bounded \term{run length} (the
longest sequence of identical symbols).
\item[Spectral shape.] The spectrum should fit the channel: little energy where the channel
loses it (high frequencies on copper), and little energy where it would radiate or couple
into other circuits.
\item[Error monitoring.] Some codes have redundancy that lets the receiver detect errors
without any higher-layer check (bipolar violations in AMI; invalid code words in 8b/10b).
\item[Framing and control.] Some codes reserve special patterns---\emph{commas}, sync
headers, control characters---that mark word boundaries and carry link-management
information.
\item[Efficiency.] All of the above costs something: extra levels, extra bandwidth or extra
bits. The \term{code rate} (information bits per transmitted bit) or the ratio of bit rate to
baud rate captures the overhead.
\end{description}

\mdcfig[0.9]{ch08_linecodes}{The classic line codes applied to the same 16-bit sequence
(dotted lines mark bit periods). NRZ codes change level at most once per bit;
Manchester codes guarantee a transition in every bit at the cost of twice the bandwidth; AMI
and MLT-3 use three levels to remove DC and, in MLT-3's case, to slow the signal down; 2B1Q
packs two bits into each four-level symbol so that the symbol rate is half the bit rate.}

\subsection{Non-return-to-zero codes}

The simplest code, \term{NRZ-L} (non-return-to-zero, level), sends $+A$ for a one and $-A$
for a zero for the full bit period. In its polar form it is the optimum binary signal in
white Gaussian noise (Section~\ref{sec:ch08:binary}), and it is spectrally compact: its
power spectral density (PSD) is $\sinc^2(fT_b)$, with a first null at the bit rate
(Figure~\ref{fig:ch08_linecode_psd}). Its problems are exactly the ones listed above: a long
run of ones or zeros produces a long constant level with no transitions and a non-zero
average. Unipolar NRZ (on--off, $0$ and $A$) is worse: it wastes half its power in a DC
component that carries no information. On-off keying of a laser or LED is unipolar NRZ, and
optical links therefore have their own DC problems, handled by the block codes and
scramblers below.

\term{NRZI} (NRZ inverted) encodes a one as a \emph{change} of level and a zero as no change
(or, in USB's convention, the reverse). The information is in the transitions, so the
receiver does not need to know the absolute polarity---a swapped wire pair does no harm. A
long run of ones now produces a square wave, but a long run of zeros still produces a
constant level, so NRZI is always paired with a code that limits the zeros: 4B5B in FDDI and
100BASE-FX, bit stuffing in USB~1.x and~2.0 (a zero is inserted after six consecutive ones),
and the (1,7) and (2,7) run-length-limited codes used on magnetic disks.

\subsection{Return-to-zero and Manchester codes}

\textbf{Return-to-zero}\index{return-to-zero} (RZ) codes return to zero in the middle of every bit. Unipolar RZ
places a half-width pulse in each one-bit and nothing in each zero-bit. Its spectrum is
twice as wide as NRZ, but it contains discrete \emph{spectral lines} at the bit rate and its
odd harmonics (Figure~\ref{fig:ch08_linecode_psd}), which a simple tuned circuit can extract
as a clock. Optical RZ formats were used in long-haul fibre systems around 2000 because
their narrow pulses were more tolerant of fibre nonlinearity.

The \term{Manchester code} (also called biphase-L) sends each bit as a pair of half-bit
levels of opposite sign: in the IEEE~802.3 convention, a one is low-then-high and a zero is
high-then-low. There is a transition in the middle of \emph{every} bit, so the code is
perfectly self-clocking and exactly DC-free whatever the data. The price is bandwidth: the
PSD, $\sinc^2(fT_b/2)\sin^2(\pi fT_b/2)$, peaks near $0.75R_b$ and has its first null at
$2R_b$. Manchester coding was used in the original 10\,Mb/s Ethernet (10BASE5, 10BASE2 and
10BASE-T), where the line ran at 20\,MBd, and it survives in RFID tags (ISO/IEC 14443 and
15693 use Manchester or related subcarrier codes), consumer infrared remotes, and the
DALI and MIL-STD-1553 buses. \textbf{Differential Manchester}\index{differential Manchester} adds the robustness of NRZI:
there is always a mid-bit transition, and a zero is marked by an extra transition at the
start of the bit. It was used in IEEE~802.5 Token Ring and survives in the AES3 (AES/EBU)
digital audio interface as \emph{biphase-mark} coding.

\begin{history}[Ethernet's Manchester choice]
When Robert Metcalfe and David Boggs built the experimental Ethernet at Xerox PARC in
1973--74, and when Digital Equipment, Intel and Xerox published the 10\,Mb/s ``DIX''
standard in 1980, clock recovery was expensive and the medium was a shared coaxial cable
with many taps. Manchester coding solved several problems at once: every station could lock
to any other station's clock from the first few bits of the preamble; the signal had no DC,
so transceivers could be transformer-isolated; and the absence of mid-bit transitions
(the ``carrier'') told a receiver that a transmission had ended. The cost---a 20\,MBd line
for 10\,Mb/s of data---was affordable on coax. It was not affordable at 100\,Mb/s on
unshielded twisted pair, which is why Fast Ethernet (1995) abandoned Manchester for 4B5B
with MLT-3.
\end{history}

\subsection{Bipolar codes: AMI, B8ZS and HDB3}

\textbf{Alternate mark inversion}\index{alternate mark inversion} (AMI), also called bipolar coding, uses three levels: a zero
is sent as no pulse, and each one (a ``mark'', in telegraph language) is sent as a pulse of
polarity opposite to the previous one. Because the ones alternate, the running sum of the
signal can never exceed one pulse, so there is no DC and very little low-frequency energy.
The symbols are correlated---a $+1$ is always followed eventually by a $-1$---and
Section~\ref{sec:ch08:psd} shows that this correlation shapes the spectrum to
$\sin^2(\pi fT_b)$, with nulls at DC and at the bit rate. The redundancy also gives free
error monitoring: two successive marks of the same polarity, a \term{bipolar violation},
can only be caused by an error. A T1 line uses AMI with 50\%-duty-cycle pulses at
1.544\,Mb/s; E1 uses the same at 2.048\,Mb/s.

AMI's weakness is a long string of zeros, which is a long silence with no transitions for the
clock recovery in the regenerative repeaters every mile or so along the line. The first T1
systems imposed a \emph{ones-density} rule on the users: no more than 15 consecutive zeros,
and at least one one in every eight bits on average, which forced 7-bit-per-channel
signalling in some cases and was the reason 56\,kb/s, not 64\,kb/s, was the usable data rate
of a DS0 for many years. Substitution codes removed the restriction:
\begin{itemize}
\item \term{B8ZS} (bipolar with eight-zero substitution, North American T1): any run of eight
zeros is replaced by the pattern \texttt{000VB0VB}, where V is a pulse that violates the
alternation rule and B is a pulse that obeys it. Two deliberate violations in a fixed pattern
cannot occur by chance from a single error, so the receiver recognises the pattern and
restores the zeros.
\item \term{HDB3} (high-density bipolar of order 3, European E1 and E3, ITU-T G.703): any run
of four zeros is replaced by \texttt{000V} or \texttt{B00V}, the choice being made so that
successive violations alternate in polarity and the DC balance is preserved. No more than
three zeros ever appear in a row.
\end{itemize}

\begin{worked}[The spectrum of a T1 line]
A T1 line carries $R_b=1.544$\,Mb/s using AMI with half-width (50\% duty) rectangular pulses.
Where are the spectral nulls and the peak, and why does it matter?

The PSD of AMI with pulse $p(t)$ is $|P(f)|^2\sin^2(\pi fT_b)/T_b$
(Section~\ref{sec:ch08:psd}). For a half-width pulse, $|P(f)|^2\propto\sinc^2(fT_b/2)$, which
has its first null at $2R_b=3.088$\,MHz. The $\sin^2(\pi fT_b)$ factor has nulls at DC and at
$R_b=1.544$\,MHz. The product is zero at DC and at 1.544\,MHz and peaks at approximately
$R_b/2=772$\,kHz. Nearly all the energy is below 1.5\,MHz.

This matters in three ways. First, with no energy at DC the line can be transformer-coupled
at both ends and can carry DC power to the repeaters on the same pairs (``simplex
powering''). Second, the energy peak at 772\,kHz sits well inside the low-loss region of
22--24-gauge cable, which at that frequency loses roughly 5--6\,dB per thousand feet; T1
repeaters were spaced about 6000\,ft apart, at the old loading-coil sites, so each span loses
about 30\,dB at 772\,kHz, and the repeater equalises and regenerates. Third, the T1 test
engineer's classic measurements of line loss and crosstalk are made at 772\,kHz, because that
is where the signal lives.
\end{worked}

\subsection{Multilevel line codes: 2B1Q, MLT-3, PAM-3 and PAM-5}

Using more than two levels reduces the symbol rate for a given bit rate, which reduces the
bandwidth and therefore the loss on a copper line.

\term{2B1Q} (two binary, one quaternary) maps each pair of bits to one of the four
levels $\pm1$ and $\pm3$; the first bit gives the sign and the second the magnitude. It is simply
4-PAM with rectangular pulses, and it was chosen for the ISDN basic-rate subscriber-loop
interface (the ``U'' interface of ANSI T1.601, 160\,kb/s at 80\,kBd) and later for HDSL. Halving the symbol rate halved the Nyquist
frequency on long subscriber loops, where loss rises steeply with frequency, and the
receiver's echo canceller (the loop is full duplex on one pair) had half as many taps to
adapt.

\term{MLT-3} (multi-level transmit, three levels) cycles through the levels
$0,+1,0,-1,0,\dots$, advancing one step for each one-bit and staying put for each zero-bit.
It takes four ones to complete a full cycle, so the fundamental frequency of the worst-case
pattern (all ones) is a quarter of the bit rate. 100BASE-TX Fast Ethernet (IEEE~802.3u,
1995) sends 125\,MBd MLT-3 over Category-5 cable, with most of the energy below about
30\,MHz (Figure~\ref{fig:ch08_linecode_psd}), which kept radiated emissions within the FCC
limits on unshielded cable. MLT-3 on its own has no guaranteed transitions (a run of zeros
is a constant level) and is not DC-free, so 100BASE-TX first maps the data through 4B5B
and then scrambles it with an 11-bit LFSR.

Three-level signalling has had a remarkable revival in the 2020s. The automotive Ethernet
standard 100BASE-T1 (IEEE~802.3bw) sends 3 bits as 2 ternary symbols (3B2T) at 66.67\,MBd
on a single twisted pair; USB4 Version~2.0 (2022) reaches 80\,Gb/s per direction using
PAM-3 with an 11-bit-to-7-trit mapping; and the JEDEC GDDR7 graphics-memory standard (2024)
uses PAM-3 as well. PAM-3 gives $\log_2 3\approx1.58$ bits per symbol, with a smaller SNR
penalty than PAM-4 and simpler receivers. Gigabit Ethernet over copper (1000BASE-T) uses
five-level PAM-5 at 125\,MBd on all four pairs at once, with a trellis code across the pairs;
10GBASE-T uses 16-level PAM at 800\,MBd with an LDPC code (Chapter~\ref{ch:wireline}).

\begin{table}[htbp]\centering\small
\caption{Classic line codes compared. $R_b$ is the bit rate. ``DC'' means whether the code
guarantees no DC component whatever the data; ``clock'' whether it guarantees transitions.}
\label{tab:ch08:linecodes}
\begin{tabularx}{\linewidth}{@{}l c c c c X@{}}
\toprule
Code & Levels & Baud rate & DC-free & Clock & Where it is (or was) used\\
\midrule
NRZ-L & 2 & $R_b$ & no & no & Inside chips, RS-232, with scrambling in SONET and SerDes\\
NRZI & 2 & $R_b$ & no & for ones & USB 1.x/2.0 (with bit stuffing), FDDI, 100BASE-FX\\
Unipolar RZ & 2 & $2R_b$ & no & for ones & Early optical links, IrDA SIR\\
Manchester & 2 & $2R_b$ & yes & yes & 10\,Mb/s Ethernet, RFID, MIL-STD-1553\\
Diff.\ Manchester & 2 & $2R_b$ & yes & yes & Token Ring, AES3 digital audio (biphase mark)\\
AMI & 3 & $R_b$ & yes & for ones & T1, E1 (with B8ZS, HDB3), ISDN S/T interface\\
2B1Q & 4 & $R_b/2$ & no & no & ISDN U-interface, HDSL\\
MLT-3 & 3 & $R_b$ & no & for ones & 100BASE-TX (after 4B5B and scrambling)\\
PAM-3 & 3 & $\approx R_b/1.5$ & -- & -- & 100BASE-T1, USB4 v2, GDDR7\\
PAM-4 & 4 & $R_b/2$ & -- & -- & 50--800G Ethernet, PCIe 6.0, GDDR6X\\
\bottomrule
\end{tabularx}
\end{table}

\subsection{Block codes: 4B5B, 8b/10b, 64b/66b and 128b/130b}

Bit-level rules like AMI and Manchester spend their redundancy inefficiently. A
\term{block line code} maps each group of $m$ data bits to a longer group of $n$ code bits,
choosing the $2^m$ code words from the $2^n$ available so that every sequence of code words
has the desired properties. The rate is $m/n$.

\subsubsection{4B5B}
The 4B5B code of FDDI (1980s) and 100BASE-TX maps each 4-bit nibble to a 5-bit code word with
at most one leading and two trailing zeros, so that with NRZI coding no more than three zeros
ever appear in a row. Of the 16 remaining 5-bit words, several serve as control symbols
(idle, start-of-stream delimiter J/K, end delimiter T/R) and the rest are invalid, which gives
error detection. The rate is 0.8: a 100\,Mb/s stream becomes 125\,MBd. 4B5B is not DC-balanced,
which is why 100BASE-TX adds a scrambler.

\subsubsection{8b/10b}
The \term{8b/10b code} of Widmer and Franaszek (IBM, 1983) is the most successful line code
ever designed. Each byte is split into a 5-bit and a 3-bit part, coded separately by a 5b/6b
and a 3b/4b table. The encoder keeps track of the \term{running disparity} (RD), the
difference between the numbers of ones and zeros sent so far. Every code word has disparity
0 or $\pm2$; for each unbalanced input there are two complementary code words, and the
encoder picks the one that drives the running disparity back towards zero. The result
(Figure~\ref{fig:ch08_rds}) is a code with
\begin{itemize}
\item strict DC balance: the running digital sum never wanders more than a few bits from zero;
\item a maximum run length of five identical bits, so that transitions are frequent and
clock recovery never starves;
\item twelve special control characters (K-codes), including the \term{comma}
K28.5 (\texttt{0011111010} or its complement), whose run of five identical bits cannot
occur across the boundary of normal data characters, so a receiver can find word
alignment by searching for it;
\item error detection: many single-bit errors produce an invalid code word or a disparity
violation.
\end{itemize}
The cost is a 25\% overhead: a 1\,Gb/s Gigabit Ethernet (1000BASE-X) link runs at 1.25\,GBd.
8b/10b appears in Fibre Channel, Gigabit Ethernet over fibre, PCI Express generations 1 and
2, Serial ATA, USB~3.0 at 5\,Gb/s, DisplayPort~1.x, JESD204B converter interfaces and
countless proprietary links. It is still the default choice for FPGA transceiver links
below a few gigabits per second.

\mdcfig{ch08_rds}{Left: the running digital sum (cumulative sum of the $\pm1$ NRZ levels)
of several bit streams. Raw random data performs a random walk; a mostly-zero idle pattern
drifts without bound, which would charge an AC-coupling capacitor; the 64b/66b scrambler
turns idle into a random walk; 8b/10b keeps the sum inside a narrow band for any data,
including the idle pattern. Right: 8b/10b limits runs to five bits, whereas a scrambled
stream has geometrically distributed runs with no hard limit.}

\begin{history}[The IBM code that ran the gigabit era]
Albert X.~Widmer and Peter A.~Franaszek, at IBM's Thomas J.~Watson Research Center, published
``A DC-balanced, partitioned-block, 8B/10B transmission code'' in the \emph{IBM Journal of
Research and Development} in 1983. Their insight was the partitioning: splitting the byte
into 5- and 3-bit sub-blocks made the encoder and decoder small enough for the logic of the
day (a few hundred gates), while careful choice of the sub-block tables kept the run length
at five and supplied the unique comma. IBM used it in the ESCON channel for mainframe
storage; it was adopted by Fibre Channel in 1994 and from there by Gigabit Ethernet in 1998.
For roughly twenty years, almost every serial link above 1\,Gb/s used it. It was finally
displaced not because it was flawed but because 25\% overhead became intolerable once lane
rates reached 10\,Gb/s.
\end{history}

\subsubsection{64b/66b and 128b/130b}
Above 10\,Gb/s, the industry switched to codes with tiny overhead that obtain their balance and
transitions statistically by scrambling. In \term{64b/66b} (IEEE~802.3ae 10GBASE-R, 2002, and
every faster Ethernet PCS since) each 64-bit block is prefixed with a 2-bit \emph{sync header},
\texttt{01} for a data block or \texttt{10} for a control block. The header is guaranteed to
contain a transition and is never scrambled, so the receiver finds block alignment by
looking for a position where every 66th bit pair is \texttt{01} or \texttt{10}. The 64
payload bits are scrambled with the self-synchronising polynomial $1+x^{39}+x^{58}$. The
overhead is 3.125\%: 10.3125\,GBd for 10\,Gb/s. PCI Express 3.0 (8\,GT/s) uses the analogous
\term{128b/130b} code with a 2-bit sync header on 128-bit blocks (1.5\% overhead) and an
additive scrambler; USB~3.1 at 10\,Gb/s uses 128b/132b. PCI Express 6.0 abandons block
coding altogether: its 64\,GT/s PAM-4 lanes carry fixed-size \emph{flits} protected by forward
error correction and CRC, with no line-code overhead at all.

A scrambled code does not \emph{guarantee} a maximum run length; it makes long runs
exponentially unlikely. With good scrambling a run of $L$ identical bits has probability of
about $2^{-L}$, so a run of 66 would be expected about once every $2^{66}$ bits---many
centuries at 10\,Gb/s. Designers therefore specify clock recovery and AC coupling to survive
runs of a few tens of bits, and check that no plausible data pattern can defeat the scrambler.

\subsection{Scramblers}
\label{sec:ch08:scramblers}

A \term{scrambler} XORs the data with a pseudo-random sequence generated by a
\term{linear-feedback shift register} (LFSR). The sequence has a period of $2^L-1$ bits if the
feedback polynomial of the $L$-stage register is \emph{primitive}, and it looks random: equal
numbers of ones and zeros, runs distributed like coin flips, and a two-valued autocorrelation.
Chapter~\ref{ch:spreadspectrum} uses the same maximal-length sequences as spreading codes.
There are two kinds of scrambler (Figure~\ref{fig:ch08_scrambler}).

\begin{figure}[htbp]\centering
\begin{tikzpicture}[scale=0.92,transform shape,
  reg/.style={draw=navy,fill=navy!6,thick,minimum height=0.6cm,minimum width=1.25cm,font=\sffamily\scriptsize},
  xor/.style={draw=accent,circle,thick,inner sep=0pt,minimum size=0.42cm,font=\small},
  arr/.style={-{Stealth[length=1.8mm]},thick,navy}]
% ----- self-synchronous scrambler
\node[font=\sffamily\small\bfseries,text=navy,anchor=west] at (-1.2,1.35) {(a) Self-synchronising (multiplicative) scrambler, $1+x^{39}+x^{58}$};
\node[xor] (x1) at (0.4,0) {$\oplus$};
\node[anchor=east,font=\sffamily\scriptsize] (din) at (-0.8,0) {data $d_n$};
\draw[arr] (din) -- (x1);
\node[reg] (r1) at (2.2,-1.1) {$s_{n-1}\cdots$};
\node[reg] (r39) at (4.2,-1.1) {$s_{n-39}$};
\node[reg] (r40) at (6.2,-1.1) {$\cdots$};
\node[reg] (r58) at (8.2,-1.1) {$s_{n-58}$};
\node[xor] (x2) at (5.2,-2.05) {$\oplus$};
\draw[arr] (x1) -- ++(1.1,0) node[right,font=\sffamily\scriptsize]{line $s_n$};
\draw[arr] (1.3,0) |- (r1.west);
\draw[arr] (r1) -- (r39); \draw[arr] (r39) -- (r40); \draw[arr] (r40) -- (r58);
\draw[arr] (r39.south) |- (x2.west);
\draw[arr] (r58.south) |- (x2.east);
\draw[arr] (x2.south) -- ++(0,-0.3) -| (x1.south);
\node[font=\sffamily\scriptsize,text=gray,anchor=west] at (9.0,-1.1) {receiver: same register};
\node[font=\sffamily\scriptsize,text=gray,anchor=west] at (9.0,-1.5) {fed by the \emph{received} bits};
% ----- additive scrambler
\node[font=\sffamily\small\bfseries,text=navy,anchor=west] at (-1.2,-3.3) {(b) Additive (synchronous) scrambler, SONET $1+x^6+x^7$, reset each frame};
\node[reg] (a1) at (1.8,-4.4) {$x^1\cdots x^5$};
\node[reg] (a6) at (3.8,-4.4) {$x^6$};
\node[reg] (a7) at (5.8,-4.4) {$x^7$};
\node[xor] (ax) at (4.8,-5.35) {$\oplus$};
\draw[arr] (a1) -- (a6); \draw[arr] (a6) -- (a7);
\draw[arr] (a6.south) |- (ax.west); \draw[arr] (a7.south) |- (ax.east);
\draw[arr] (ax.south) -- ++(0,-0.3) -| (0.7,-4.4) -- (a1.west);
\node[xor] (ay) at (8.3,-4.4) {$\oplus$};
\draw[arr] (a7.east) -- (ay.west);
\node[anchor=west,font=\sffamily\scriptsize] (d2) at (8.8,-3.75) {data};
\draw[arr] (d2.west) -| (ay.north);
\draw[arr] (ay.east) -- ++(1.0,0) node[right,font=\sffamily\scriptsize]{line};
\end{tikzpicture}
\caption{(a) A self-synchronising scrambler feeds its own output back through the
register; the descrambler is the mirror image fed by the received bits and needs no
initialisation, but each line error is multiplied (here into three errors). (b) An
additive scrambler generates a free-running pseudo-random sequence and XORs it with the
data; there is no error multiplication, but transmitter and receiver registers must be
synchronised, for example by resetting them at a known point in each frame.}
\label{fig:ch08_scrambler}
\end{figure}

A \term{self-synchronising scrambler} (or multiplicative scrambler) computes
$s_n=d_n\oplus s_{n-39}\oplus s_{n-58}$ in the 64b/66b case: the line bit depends on the data
bit and on earlier \emph{line} bits. The descrambler computes
$d_n=s_n\oplus s_{n-39}\oplus s_{n-58}$ from the received line bits, so after 58 bits it is
automatically in step with the scrambler, whatever state either started in. This is
convenient---there is nothing to synchronise---but it has two costs. Each line error appears
three times at the descrambler output (at $n$, $n+39$ and $n+58$), which the CRC or FEC above
must tolerate; Ethernet's CRC-32 detects any three bit errors in a maximum-length frame, which
is part of why this polynomial was acceptable. Its great advantage is that, with 58 stages,
no practical data pattern can accidentally match the scrambler state for long.

An \term{additive scrambler} (synchronous scrambler) runs its LFSR independently of the data
and simply XORs its output onto the stream. There is no error multiplication, but the two
registers must be synchronised. SONET/SDH uses the 7-stage generator $1+x^6+x^7$ (period 127),
reset to all ones at the start of every frame. Its period is so short that users could, and
did, send packet payloads that reproduced the scrambler sequence and produced long runs of
zeros---the ``killer packet'' problem---which is why Packet-over-SONET (RFC~2615) added a
second, self-synchronising $1+x^{43}$ scrambler on the payload. PCI Express~1 and~2 use $1+x^3+x^4+x^5+x^{16}$,
reset by special symbols; PCI Express~3 to~5 use the 23-stage polynomial
$1+x^2+x^5+x^8+x^{16}+x^{21}+x^{23}$ with a different seed per lane so that adjacent lanes do
not radiate identical patterns. DVB-S2 baseband frames use an additive $1+x^{14}+x^{15}$
scrambler, and 802.11 Wi-Fi uses $1+x^4+x^7$ with a random seed per packet.

\begin{pitfall}[A scrambler is not a line code]
Scrambling makes data look random; it does not make it DC-free or run-limited. A short
scrambler cannot prevent a user from sending data that happens to equal its sequence, as the
SONET killer-packet episode showed. Robust designs use long registers, scramble the payload
with a self-synchronising scrambler that no fixed data pattern can track, and use receivers
that tolerate runs of tens of bits. When you design a custom FPGA link, test it with
worst-case patterns, not just PRBS-31.
\end{pitfall}

\begin{inpractice}[PRBS test patterns]
The same LFSRs generate the \term{pseudo-random binary sequences} (PRBS) that every bit-error
tester, oscilloscope and SerDes built-in self-test uses: PRBS-7 ($1+x^6+x^7$), PRBS-15
($1+x^{14}+x^{15}$), PRBS-23 ($1+x^{18}+x^{23}$) and PRBS-31 ($1+x^{28}+x^{31}$), as defined in
ITU-T O.150 and the relevant IEEE and OIF documents. The longer the pattern, the longer its
longest run and the richer its low-frequency content, so PRBS-31 is the stress test for
baseline wander and clock recovery, and PRBS-7 the quick check. For PAM-4 links, the
standards define PRBS13Q and PRBS31Q (Gray-coded PAM-4 versions) and the SSPRQ ``short
stress pattern'' used for transmitter compliance.
\end{inpractice}

%======================================================================
\section{The power spectrum of a PAM signal}
\label{sec:ch08:psd}
%======================================================================

Every claim about line-code spectra in the previous section follows from one formula. It is
worth deriving carefully, because it also explains spectral lines, the spectral effect of
correlation between symbols, and later the spectra of QAM and OFDM signals.

\subsection{Cyclostationarity and the averaged autocorrelation}

The PAM signal $s(t)=\sum_k a_kp(t-kT)$ with random symbols is not wide-sense stationary: its
statistics repeat with period $T$ (the variance is higher at the pulse centres than between
them, for example). Such a process is \term{cyclostationary}. Its PSD is defined through the
autocorrelation averaged over one period, which is what a spectrum analyser measures. Assume
the symbols form a wide-sense stationary sequence with mean $\mu_a=\E[a_k]$ and
autocorrelation $R_a[m]=\E[a_k a_{k+m}]$ (real symbols for simplicity). Then
\begin{align}
R_s(t+\tau,t)&=\E[s(t+\tau)s(t)]=\sum_k\sum_l R_a[l-k]\,p(t+\tau-lT)\,p(t-kT),\notag\\
\bar R_s(\tau)&=\frac1T\int_0^T R_s(t+\tau,t)\,dt
=\frac1T\sum_m R_a[m]\int_{-\infty}^{\infty}p(u+\tau-mT)\,p(u)\,du,
\end{align}
where the second line uses the fact that summing $k$ over all integers and integrating $t$ over
one period is the same as integrating over all time. The integral is the pulse's deterministic
autocorrelation $r_p(\tau-mT)$, whose Fourier transform is $|P(f)|^2$. Taking the transform
gives the \term{PAM spectrum formula}
\begin{equation}
\boxed{\;S_s(f)=\frac{|P(f)|^2}{T}\,S_a(f),\qquad S_a(f)=\sum_{m=-\infty}^{\infty}R_a[m]\,e^{-j2\pi fmT}.\;}
\label{eq:ch08:psd}
\end{equation}
The spectrum factors into two parts: the \emph{pulse} term $|P(f)|^2/T$ and the \emph{symbol}
term $S_a(f)$, the discrete-time PSD of the symbol sequence (periodic in $f$ with period
$1/T$). The pulse sets the envelope; the correlation between symbols sculpts it.

\subsection{Uncorrelated symbols and spectral lines}

For independent symbols with variance $\sigma_a^2$ and mean $\mu_a$, $R_a[m]=\sigma_a^2\delta[m]+\mu_a^2$.
The constant $\mu_a^2$ term is an infinite sum of complex exponentials, which by the Poisson
sum formula is a train of impulses: $\sum_m e^{-j2\pi fmT}=\frac1T\sum_k\delta(f-k/T)$.
Substituting into~\eqref{eq:ch08:psd},
\begin{equation}
S_s(f)=\frac{\sigma_a^2}{T}|P(f)|^2+\frac{\mu_a^2}{T^2}\sum_{k=-\infty}^{\infty}\big|P(k/T)\big|^2\,\delta\!\left(f-\frac kT\right).
\label{eq:ch08:psdlines}
\end{equation}
This is the most useful result about line-code spectra. The first term is a continuous
spectrum shaped like the pulse. The second term exists only if the symbols have a non-zero
mean: it is a set of \term{spectral lines} at multiples of the symbol rate, weighted by the
pulse spectrum at those frequencies. Several consequences follow at once
(Figure~\ref{fig:ch08_linecode_psd}):
\begin{itemize}
\item \emph{Polar} codes with equiprobable symbols have $\mu_a=0$ and no lines. Polar NRZ has
$S(f)=A^2T\sinc^2(fT)$, with most of its power below the bit rate (about 90\% inside
$|f|<1/T$).
\item \emph{Unipolar NRZ} ($a\in\{0,A\}$) has $\sigma_a^2=\mu_a^2=A^2/4$. Its continuous part is
a quarter of polar NRZ's, and it has a DC line of power $A^2/4$---half the total power
spent on a constant that carries no information. The other lines fall on the nulls of
$\sinc^2(fT)$ and vanish.
\item \emph{Unipolar RZ} with half-width pulses has $|P(f)|^2=(AT/2)^2\sinc^2(fT/2)$, whose
nulls are at even multiples of $1/T$. The lines at $f=\pm1/T,\pm3/T,\dots$ survive. The line
at the bit rate has power $A^2/(4\pi^2)$ in each sideband, a ready-made clock tone that a
tuned circuit or narrowband PLL can extract directly.
\end{itemize}

\mdcfig{ch08_linecode_psd}{Power spectral densities of common line codes (unit amplitude,
equiprobable bits; lines are theory from~\eqref{eq:ch08:psd}--\eqref{eq:ch08:psdlines} and
dots are Welch estimates from $2\times10^5$ simulated bits). Left: polar codes. Manchester
moves the energy up and away from DC; AMI and MLT-3 use symbol correlation to create a null at
DC, with MLT-3's energy concentrated below about $0.3R_b$. Right: unipolar codes waste power
in spectral lines (arrows); the RZ line at the bit rate is a free clock tone.}

\begin{keyidea}[Spectral lines come from the mean; spectral nulls come from correlation]
A non-zero symbol mean produces discrete lines at multiples of the symbol rate, weighted by
$|P(k/T)|^2$. They cost power but can supply a clock. Correlation between symbols multiplies
the continuous spectrum by $S_a(f)$, which can place nulls where the channel cannot pass
energy---at DC for a transformer-coupled line, or at the band edge for a partial-response
channel. Designing a line code is designing $S_a(f)$.
\end{keyidea}

\subsection{Correlated symbols: AMI, Manchester and MLT-3}

For AMI, the symbols take the values $0$ (probability $\tfrac12$) and $\pm1$ (probability
$\tfrac14$ each), and successive marks alternate. A short calculation gives $R_a[0]=\tfrac12$,
$R_a[\pm1]=-\tfrac14$ (the product $a_ka_{k+1}$ is $-1$ when both bits are ones, with
probability $\frac14$, and zero otherwise), and $R_a[m]=0$ for $|m|\ge2$ (a pair of marks
separated by one or more bits has equally likely equal and opposite polarities). Hence
\begin{equation}
S_a(f)=\tfrac12-\tfrac12\cos(2\pi fT)=\sin^2(\pi fT),
\end{equation}
which has a double zero at DC. AMI with NRZ-width pulses has
$S(f)=T\sinc^2(fT)\sin^2(\pi fT)$; the negative correlation between neighbours has pushed the
energy away from DC without widening the spectrum.

Manchester can be analysed either as a binary code with the two-half-bit pulse
$p(t)=\mathrm{rect}$ on the second half minus $\mathrm{rect}$ on the first, which gives
$|P(f)|^2=T^2\sinc^2(fT/2)\sin^2(\pi fT/2)$ and hence
$S(f)=T\sinc^2(fT/2)\sin^2(\pi fT/2)$, or as a correlated code at twice the rate. The first
view is simpler: it is the \emph{pulse} that has no DC content, because its area is zero, so
no data pattern can produce DC.

MLT-3 is a Markov code: its state walks around the cycle $0,+1,0,-1$ one step per one-bit.
Writing the level as $\cos(\pi\theta_k/2)$ with the state $\theta_k$ uniform on $\{0,1,2,3\}$
and averaging over the random walk gives
$R_a[m]=\tfrac12\,\re\!\big[\big(\tfrac{1+j}{2}\big)^{|m|}\big]$: a damped cosine at a quarter
of the bit rate. The resulting spectrum (Figure~\ref{fig:ch08_linecode_psd}) peaks at about
$0.11R_b$ and is small above $0.3R_b$. For 100BASE-TX at 125\,MBd that means a peak near
14\,MHz and little energy above about 40\,MHz, which is why Category-5 cable rated to
100\,MHz was enough.

\begin{pitfall}[A spectrum analyser measures something slightly different]
The averaged PSD of~\eqref{eq:ch08:psd} assumes random data. A real link idling on a fixed
pattern (a comma, an idle ordered set, a repeated training sequence) is periodic, so its
spectrum is a set of lines, and those lines can fail an electromagnetic-emissions test that
random data would pass. This is one reason every modern link scrambles its idle pattern, and
why spread-spectrum clocking (a slow frequency modulation of the reference clock by about
0.5\%) is used in PCI Express, SATA and DisplayPort to smear the remaining lines.
\end{pitfall}

%======================================================================
\section{Intersymbol interference and the Nyquist criterion}
\label{sec:ch08:nyquist}
%======================================================================

\subsection{What ISI looks like}

A rectangular pulse of duration $T$ is perfect in time---neighbouring pulses do not overlap---but
its $\sinc$ spectrum extends to infinity. Any real channel is band-limited, and band-limiting a
pulse spreads it in time. In Figure~\ref{fig:ch08_isi} a train of rectangular pulses passes
through a first-order RC low-pass with a time constant of $0.6T$, a crude model of a cable.
Each pulse rises too slowly to reach full amplitude and leaves a tail that decays into the next
symbol periods. The sample of each symbol is now a sum of that symbol and the tails of its
predecessors, exactly the middle term of~\eqref{eq:ch08:sample}. Whether a given symbol is
decided correctly depends on the data around it. When the tails are large enough, some data
patterns cause errors even without noise.

\mdcfig{ch08_isi}{Left: $\sinc$ pulses overlap heavily, yet at each sampling instant every
pulse except one passes through zero, so the samples (red) equal the transmitted symbols
exactly. Right: rectangular pulses smeared by an RC low-pass leave tails that add to the
following samples; the dotted lines show the error each sample suffers from the tails of its
predecessors.}

The left half of Figure~\ref{fig:ch08_isi} shows the way out. Pulses may overlap as much as
they like in time, provided each pulse passes through \emph{zero} at every sampling instant of
the other symbols. Then the sampled sum is exactly the wanted symbol. This is the idea behind
Nyquist's 1928 paper.

\subsection{The Nyquist criterion}

Let $x(t)$ be the end-to-end pulse of~\eqref{eq:ch08:overall}, with the sampling instant taken
as $t_0=0$ for convenience. There is no ISI if and only if
\begin{equation}
x(kT)=\begin{cases}1,&k=0,\\0,&k\ne0.\end{cases}
\label{eq:ch08:nyqtime}
\end{equation}
To translate this into the frequency domain, sample $x(t)$ with an impulse train. Sampling in
time replicates the spectrum in frequency (Chapter~\ref{ch:pcm}), so the samples $x(kT)$ have
the discrete-time Fourier transform $\frac1T\sum_k X(f-k/T)$. The condition
$x(kT)=\delta[k]$ says this transform is the constant 1. Therefore:

\begin{keyidea}[The Nyquist criterion for zero ISI]
A pulse $x(t)$ with $x(0)=1$ has zero crossings at every non-zero multiple of $T$ if and only if
its spectrum satisfies
\begin{equation}
\sum_{k=-\infty}^{\infty}X\!\left(f-\frac kT\right)=T\qquad\text{for all }f.
\label{eq:ch08:nyquist}
\end{equation}
The shifted copies of the spectrum must add up to a flat line. A pulse that satisfies this is
called a \term{Nyquist pulse}.
\end{keyidea}

Two consequences are immediate.
\begin{enumerate}
\item \emph{Minimum bandwidth.} If $X(f)$ is confined to $|f|<W$ with $W<1/(2T)$, the copies do
not touch and there are gaps where the sum is zero, so~\eqref{eq:ch08:nyquist} fails. A pulse
that carries $R_s=1/T$ symbols per second without ISI needs a bandwidth of at least
\begin{equation}
W_{\min}=\frac{1}{2T}=\frac{R_s}{2},
\end{equation}
the \term{Nyquist bandwidth}. Equivalently, a channel of bandwidth $W$ can carry at most $2W$
independent symbols per second without ISI---the \emph{Nyquist rate}, the same $2W$ as in the
sampling theorem, and for the same reason.
\item \emph{The minimum-bandwidth pulse is unique.} At exactly $W=1/(2T)$ the only solution is
the brick wall $X(f)=T$ for $|f|<1/(2T)$, whose impulse response is $x(t)=\sinc(t/T)$. It is
impractical: its tails decay only as $1/t$, so a truncated version needs a very long filter, and
the sum of the tails of many symbols can diverge if the sampling instant is even slightly wrong
(Section~\ref{sec:ch08:rc}).
\end{enumerate}
With more bandwidth than the minimum there is freedom: any spectrum whose roll-off region
around $1/(2T)$ has \emph{odd symmetry}---whatever it loses below $1/(2T)$ it gains back as a
mirror image above---satisfies~\eqref{eq:ch08:nyquist}. This is the \term{vestigial symmetry}
condition, and it is the key to practical pulses.

\begin{history}[Nyquist, 1928]
Harry Nyquist, a Swedish-born engineer at AT\&T and later Bell Labs, published ``Certain
factors affecting telegraph speed'' in the \emph{Bell System Technical Journal} in 1924 and
``Certain topics in telegraph transmission theory'' in the \emph{Transactions of the AIEE} in
1928. The second paper showed that $2W$ independent pulse values per second can be sent through
a channel of bandwidth $W$, and gave the vestigial-symmetry condition for pulses that do so
without interference. Nyquist was thinking about telegraph signals on cables, not about
sampling, but his result is the direct ancestor of the sampling theorem and, via Hartley's
1928 paper on information, of Shannon's capacity formula. The same Nyquist also explained
thermal noise (with John B.~Johnson, 1928) and invented the stability criterion used in every
feedback-control course (1932). Few engineers have their name attached to three different
fundamental results.
\end{history}

%======================================================================
\section{Raised-cosine and root-raised-cosine pulses}
\label{sec:ch08:rc}
%======================================================================

\subsection{The raised-cosine family}

The standard Nyquist pulse is the \term{raised cosine} (RC). Its spectrum is flat up to
$(1-\beta)/(2T)$, falls as a half-period of a cosine through the Nyquist frequency $1/(2T)$,
and reaches zero at $(1+\beta)/(2T)$:
\begin{equation}
X_{\mathrm{RC}}(f)=\begin{cases}
T, & |f|\le\dfrac{1-\beta}{2T},\\[2mm]
\dfrac T2\left[1+\cos\!\left(\dfrac{\pi T}{\beta}\left(|f|-\dfrac{1-\beta}{2T}\right)\right)\right], & \dfrac{1-\beta}{2T}<|f|\le\dfrac{1+\beta}{2T},\\[2mm]
0, & \text{otherwise.}
\end{cases}
\label{eq:ch08:rcspec}
\end{equation}
The parameter $\beta\in[0,1]$ is the \term{roll-off factor} or \term{excess bandwidth}: the
occupied bandwidth is
\begin{equation}
B=\frac{1+\beta}{2T}=\frac{R_s(1+\beta)}{2}\ \ \text{(baseband)},\qquad
B=R_s(1+\beta)\ \ \text{(passband, on a carrier).}
\label{eq:ch08:rcbw}
\end{equation}
The impulse response is
\begin{equation}
x_{\mathrm{RC}}(t)=\sinc\!\left(\frac tT\right)\frac{\cos(\pi\beta t/T)}{1-(2\beta t/T)^2}.
\label{eq:ch08:rctime}
\end{equation}
The $\sinc(t/T)$ factor supplies the zero crossings at multiples of $T$; the second factor
tapers the tails, which now decay as $1/|t|^3$ instead of $1/|t|$
(Figure~\ref{fig:ch08_raised_cosine}). The case $\beta=0$ is the brick-wall sinc; $\beta=1$ is
the ``full'' raised cosine, which has a spectrum twice the Nyquist bandwidth and additional
zero crossings halfway between the sampling instants.

\mdcfig{ch08_raised_cosine}{The raised-cosine family. Left: impulse responses; all pass
through zero at every non-zero multiple of $T$. Centre: spectra; the roll-off is odd-symmetric
about the Nyquist frequency $1/(2T)$, so the shifted copies sum to a constant. Right: the tail
envelope on a log scale; larger roll-off buys much faster decay.}

\subsection{Choosing the roll-off}

The roll-off is one of the most consequential numbers in a physical-layer specification. A
small $\beta$ saves bandwidth; a large $\beta$ buys:
\begin{itemize}
\item \emph{Shorter filters.} The tails decay faster, so the pulse can be truncated to fewer
symbol periods for a given level of residual ISI and out-of-band emission.
\item \emph{Timing tolerance.} With a small roll-off, the tails are large, and if the receiver
samples a fraction of a symbol away from the ideal instant, the tails of many neighbouring
symbols no longer pass through zero and add up. Figure~\ref{fig:ch08_timing_sensitivity}
shows the worst-case eye opening against timing error: at $\beta=0.1$ a timing error of 0.2
symbols closes the eye completely, whereas at $\beta=1$ the eye loses only a few decibels.
\item \emph{Lower peak-to-average power ratio.} Large tails also mean large peaks when many of
them add in phase. In a simulation of RRC-shaped QPSK, the peak-to-average power ratio exceeded
with probability $10^{-4}$ is about 3.7\,dB at $\beta=0.35$, 5.1\,dB at $\beta=0.20$ and
6.3\,dB at $\beta=0.05$; every extra decibel costs power-amplifier back-off
(Chapter~\ref{ch:sdr}).
\end{itemize}

\mdcfig[0.7]{ch08_timing_sensitivity}{Worst-case eye opening (the smallest value of the
decision variable over a long random sequence, relative to the ideal) against sampling-time
error, for raised-cosine pulses. Small roll-off demands precise timing.}

\begin{inpractice}[Roll-off factors in real systems]
Wideband CDMA (3GPP TS 25.104) uses $\beta=0.22$ at 3.84\,Mchip/s, giving an occupied
bandwidth of 4.68\,MHz inside its 5\,MHz channel. TETRA uses 0.35. DVB-S (1994) used 0.35;
DVB-S2 (ETSI EN 302 307, 2005) offers 0.35, 0.25 and 0.20; and DVB-S2X (2014) adds 0.15, 0.10
and 0.05, because on a satellite transponder of fixed bandwidth, lower roll-off translates
directly into more symbols per second. Cable-TV QAM (ITU-T J.83 Annex~B) uses 0.18 for 64-QAM
and 0.12 for 256-QAM. The trend is towards smaller roll-offs as digital filters become cheap
and receivers' timing recovery improves. At the other extreme, high-speed wireline links
rarely use explicit raised-cosine filtering at all: the channel itself provides the
band-limiting and an equaliser shapes the result (Section~\ref{sec:ch08:channel}).
\end{inpractice}

\begin{worked}[Roll-off and a satellite transponder]
A 36\,MHz satellite transponder carries a single DVB-S2 carrier. What symbol rate fits with
$\beta=0.35$, 0.20 and 0.05, and how much extra throughput does the lowest roll-off give?

From~\eqref{eq:ch08:rcbw}, the passband bandwidth is $R_s(1+\beta)$, so $R_s=36/(1+\beta)$\,MBd:
26.7\,MBd at $\beta=0.35$, 30.0\,MBd at 0.20 and 34.3\,MBd at 0.05. In practice a guard
band of a few per cent is left between carriers, but the ratio stands: going from 0.35 to
0.05 gives $1.35/1.05=1.29$, a 29\% throughput increase for the same modulation and coding,
and 14\% compared with 0.20. The price is a filter roughly three to four times longer, tighter
timing recovery, and more amplifier back-off. For a satellite operator who charges by the
megahertz, that is a bargain.
\end{worked}

\subsection{Splitting the filter: the root raised cosine}

The Nyquist criterion constrains the \emph{end-to-end} pulse $x(t)=g_T*c*g_R$. The next section
shows that, in white noise, the receive filter should be matched to the received pulse,
$G_R(f)=G_T^*(f)$ when the channel is flat. Both requirements are met by splitting the
raised-cosine spectrum equally:
\begin{equation}
|G_T(f)|=|G_R(f)|=\sqrt{X_{\mathrm{RC}}(f)}.
\end{equation}
This is the \term{root raised cosine} (RRC), used in nearly every single-carrier radio
standard. Its impulse response is
\begin{equation}
g_{\mathrm{RRC}}(t)=\frac{1}{\sqrt T}\,
\frac{\sin\!\big(\pi(1-\beta)t/T\big)+4\beta(t/T)\cos\!\big(\pi(1+\beta)t/T\big)}
{\pi(t/T)\big[1-(4\beta t/T)^2\big]},
\label{eq:ch08:rrc}
\end{equation}
with the limiting values at $t=0$ and $t=\pm T/(4\beta)$ found by l'H\^opital's rule. The RRC
pulse alone is \emph{not} a Nyquist pulse---it has ISI---but the cascade of two RRC filters is.
A receiver that looks at the transmitted waveform before the matched filter will therefore see
a messy eye; only after the receive RRC does the eye open. Figure~\ref{fig:ch08_eyes} shows
the result for three roll-offs.

\mdcfig{ch08_eyes}{Eye diagrams after an RRC transmit filter and RRC matched filter
(truncated to 16 symbols), for three roll-off factors. Top: without noise the eye is fully
open at the sampling instant, with thin crossings for large $\beta$. Bottom: at
$E_s/N_0=15$\,dB, the same noise gives the same vertical spread, but the small-roll-off eyes
close faster away from the centre.}

\begin{pitfall}[Truncation and quantisation break the Nyquist property]
A digital RRC filter is a truncated, sampled and quantised version of~\eqref{eq:ch08:rrc}.
Truncating it leaves residual ISI and raises the stopband: with $\beta=0.22$ and plain
truncation, a span of 8 symbols leaves sidelobes only about 30\,dB down, and even 20 symbols
reaches only about 48\,dB. Designers often window the truncated response or optimise the taps
directly (for example with a least-squares design against the ideal spectrum) and then check
two things: the adjacent-channel emission against the spectrum mask, and the residual ISI of
the TX--RX cascade, usually quoted as the \emph{error-vector magnitude} (EVM) with no noise.
A self-ISI EVM below about 1\% ($-40$\,dB) is a typical target. Remember also that the
transmitter's analog reconstruction filter and the receiver's anti-aliasing filter are part of
the cascade and have phase distortion of their own.
\end{pitfall}

%======================================================================
\section{The matched filter}
\label{sec:ch08:mf}
%======================================================================

The Nyquist criterion removes ISI. The receive filter has a second job: to let through as
much signal and as little noise as possible. Surprisingly, the answer to ``what is the best
receive filter?'' does not depend on the details of the noise level or the alphabet. It depends
only on the pulse.

\subsection{Derivation by the Cauchy--Schwarz inequality}

Consider one pulse $p(t)$ of energy $E_p=\int|p(t)|^2dt$, received in additive white Gaussian
noise of two-sided PSD $N_0/2$. The receiver passes $r(t)=p(t)+n(t)$ through a filter $h(t)$
and samples the output at time $T_0$. The signal part of the sample is
\begin{equation}
y_s(T_0)=\int_{-\infty}^{\infty}h(\tau)\,p(T_0-\tau)\,d\tau,
\end{equation}
and the noise part has variance
\begin{equation}
\sigma^2=\frac{N_0}{2}\int_{-\infty}^{\infty}|h(\tau)|^2\,d\tau.
\end{equation}
The output signal-to-noise ratio is therefore
\begin{equation}
\mathrm{SNR}=\frac{\big|\int h(\tau)\,p(T_0-\tau)\,d\tau\big|^2}{\frac{N_0}{2}\int|h(\tau)|^2d\tau}.
\end{equation}
The \term{Cauchy--Schwarz inequality} says that for any two finite-energy functions,
$\big|\int u\,v\,d\tau\big|^2\le\int|u|^2d\tau\int|v|^2d\tau$, with equality if and only if one is
a scalar multiple of the conjugate of the other. Applying it with $u=h(\tau)$ and
$v=p(T_0-\tau)$,
\begin{equation}
\mathrm{SNR}\le\frac{\int|h|^2d\tau\;\int|p(T_0-\tau)|^2d\tau}{\frac{N_0}{2}\int|h|^2d\tau}=\frac{2E_p}{N_0},
\label{eq:ch08:mfsnr}
\end{equation}
with equality when
\begin{equation}
\boxed{\;h(t)=K\,p^*(T_0-t).\;}
\label{eq:ch08:mf}
\end{equation}

\begin{keyidea}[The matched filter]
In white Gaussian noise, the linear filter that maximises the sampled signal-to-noise ratio
has an impulse response equal to the time-reversed, conjugated, delayed pulse. Its output SNR
is $2E_p/N_0$, which depends only on the pulse \emph{energy}, not on its shape, bandwidth or
duration. In the frequency domain, $H(f)=K\,P^*(f)e^{-j2\pi fT_0}$: the filter emphasises the
frequencies where the signal is strong and ignores those where there is only noise.
\end{keyidea}

This result is remarkable. Two pulses of the same energy---a short, tall one and a long, flat
one; a chirp and a rectangle---are equally detectable in white noise if each receiver is
matched to its pulse. That is why radar can send long, low-power chirps instead of short,
high-power spikes (pulse compression), why GPS signals 20\,dB below the thermal noise floor can
be received (Chapter~\ref{ch:spreadspectrum}), and why deep-space probes trade time for
energy at rates of a few bits per second. Matched filtering is the reason energy, and hence
$E_b/N_0$, is the natural currency of digital communications.

\subsection{The correlator and integrate-and-dump}

Sampling the matched filter's output at $T_0$ gives
\begin{equation}
y(T_0)=\int r(\tau)\,h(T_0-\tau)\,d\tau=\int r(\tau)\,p^*(\tau)\,d\tau,
\end{equation}
which is simply the \term{correlation} of the received signal with the known pulse. A
\term{correlator receiver} multiplies $r(t)$ by a locally generated replica of $p(t)$ and
integrates. For a rectangular pulse the replica is a constant and the correlator reduces to
an \term{integrate-and-dump} circuit: integrate the received signal over each bit, sample,
reset. Figure~\ref{fig:ch08_matched_filter} shows it at work on a signal whose per-sample SNR is
about $-1$\,dB---individual samples are almost useless---but whose integrated value at the end of
each bit is decided correctly every time. Integrating 32 independent noise samples improves the
SNR by $10\log_{10}32=15$\,dB.

\mdcfig{ch08_matched_filter}{The matched filter for a rectangular pulse is an
integrate-and-dump. Top: polar NRZ at a per-sample SNR of about $-1$\,dB. Bottom: the
integrator output ramps towards $\pm A$ during each bit and is sampled (red dots) at the end of
the bit, then reset. The noise averages out; the signal adds coherently.}

For a filter matched to a pulse that is also a Nyquist pulse after matching---the RRC pair of
Section~\ref{sec:ch08:rc}---the matched filter maximises SNR \emph{and} gives zero ISI. That is
the condition for a pulse $g(t)$ whose autocorrelation $\int g(\tau)g^*(\tau-kT)d\tau$ is zero
for $k\ne0$: in other words, the shifted pulses $g(t-kT)$ are \emph{orthogonal}. The RRC pulse
is the standard orthogonal pulse, and a PAM signal built from it is an orthonormal expansion
whose coefficients the matched filter recovers exactly, sample by sample. This geometric view
is developed in Chapter~\ref{ch:modulation}.

\subsection{Coloured noise and the whitening filter}

If the noise is not white but has PSD $S_n(f)$, the same argument applied in the frequency
domain (Cauchy--Schwarz on $\int H P\,e^{j2\pi fT_0}df$ with weights $S_n(f)$) gives
\begin{equation}
H(f)=K\,\frac{P^*(f)}{S_n(f)}\,e^{-j2\pi fT_0},\qquad
\mathrm{SNR}_{\max}=\int_{-\infty}^{\infty}\frac{|P(f)|^2}{S_n(f)}\,df.
\label{eq:ch08:coloured}
\end{equation}
This filter factors into two stages. A \term{whitening filter} $1/\sqrt{S_n(f)}$ makes the noise
white and changes the pulse to $P(f)/\sqrt{S_n(f)}$; a filter matched to that whitened pulse
follows. The result is intuitive: frequencies where the noise is strong are attenuated, even if
the signal is strong there too. It matters in practice whenever crosstalk or interference
dominates thermal noise---on a DSL line, where near-end crosstalk from other pairs in the same
binder rises with frequency, or in a receiver with a strong narrowband interferer. The
same structure---whiten, then match---appears in the equalisers of Chapter~\ref{ch:equalization}
and in the optimum MIMO receivers of Chapter~\ref{ch:mimo}.

\begin{pitfall}[Matched to what?]
The matched filter must be matched to the pulse \emph{as received}, after the channel. In a
radio with a flat channel the receive RRC is matched to the transmit RRC and all is well. On a
cable with 20\,dB of high-frequency loss, a filter matched to the \emph{transmitted} pulse is
badly mismatched: the received pulse is long and smeared. The combination of a filter matched
to the received pulse followed by a symbol-spaced equaliser is optimal; that structure, the
\emph{whitened matched filter}, is the front end of the maximum-likelihood sequence detector of
Chapter~\ref{ch:equalization}.
\end{pitfall}

%======================================================================
\section{Binary detection in Gaussian noise}
\label{sec:ch08:binary}
%======================================================================

After the matched filter and sampler, the receiver has one number per symbol, $y=a+\nu$, where
$a$ is the noiseless sample value and $\nu$ is Gaussian with variance $\sigma^2$. The last step
is to decide which symbol was sent.

\subsection{MAP and ML decisions}

Suppose one of two values $a_0<a_1$ is sent with prior probabilities $P_0$ and $P_1$. The
decision that minimises the probability of error is the \term{maximum a posteriori} (MAP) rule:
choose the symbol with the larger posterior probability $P_i\,f(y\,|\,a_i)$. With Gaussian
likelihoods $f(y|a_i)\propto e^{-(y-a_i)^2/2\sigma^2}$, taking logarithms reduces the comparison to
a threshold test, $y\gtrless\gamma$, with
\begin{equation}
\gamma=\frac{a_0+a_1}{2}+\frac{\sigma^2}{a_1-a_0}\ln\frac{P_0}{P_1}.
\label{eq:ch08:map}
\end{equation}
With equal priors the log term vanishes and the threshold is the midpoint---the
\term{maximum-likelihood} (ML) rule, used unless the source is known to be biased. Unequal
priors shift the threshold towards the less likely symbol, by an amount that grows with the
noise. (Modern soft-decision decoders, Chapter~\ref{ch:moderncodes}, do not make this hard
decision at all; they pass the log-likelihood ratio $\ln\frac{f(y|a_1)}{f(y|a_0)}=\frac{(a_1-a_0)}{\sigma^2}\big(y-\frac{a_0+a_1}2\big)$ to the decoder.)

With the midpoint threshold, an error occurs when the noise pushes $y$ past the midpoint, a
distance $d/2$ from each point, where $d=a_1-a_0$. The probability is
\begin{equation}
P_e=\Q\!\left(\frac{d}{2\sigma}\right),\qquad \Q(x)=\frac{1}{\sqrt{2\pi}}\int_x^\infty e^{-u^2/2}du.
\label{eq:ch08:pe}
\end{equation}
The \term{Q-function}, the tail of the standard Gaussian, is the most important function in
this book (Chapter~\ref{ch:noise}). Values worth memorising: $\Q(x)=10^{-3}$ at $x=3.09$,
$10^{-6}$ at 4.75, $10^{-9}$ at 6.00, $10^{-12}$ at 7.03 and $10^{-15}$ at 7.94. For large
$x$, $\Q(x)\approx e^{-x^2/2}/(x\sqrt{2\pi})$, so the error probability falls exponentially
with the square of the distance-to-noise ratio: each additional decibel of SNR buys more and
more decades of error rate.

\subsection{Energy bookkeeping: \texorpdfstring{$E_b/N_0$}{Eb/N0}}

To compare signalling schemes fairly, express $d/2\sigma$ in terms of the energy spent per bit.
With a matched filter normalised to unit energy, the noise variance at the sampler is
$\sigma^2=N_0/2$ and a pulse of energy $E$ produces a sample of amplitude $\sqrt E$. For three
classic binary schemes:
\begin{itemize}
\item \emph{Antipodal} (polar NRZ, BPSK): the two pulses are $\pm p(t)$, each of energy $E_b$.
The samples are $\pm\sqrt{E_b}$, $d=2\sqrt{E_b}$, and
\begin{equation}
P_b=\Q\!\left(\sqrt{\frac{2E_b}{N_0}}\right).
\label{eq:ch08:antipodal}
\end{equation}
\item \emph{Orthogonal} (e.g.\ two tones in coherent FSK, or pulses in different time slots as
in pulse-position modulation): the pulses $p_0(t)$, $p_1(t)$ each have energy $E_b$ and zero
cross-correlation, so they are a distance $\sqrt{2E_b}$ apart in signal space
(Chapter~\ref{ch:modulation}), and
\begin{equation}
P_b=\Q\!\left(\sqrt{\frac{E_b}{N_0}}\right).
\end{equation}
\item \emph{On--off} (unipolar NRZ, on-off keying): the pulses are $0$ and $p(t)$ with peak
energy $E_1$. The average energy is $E_b=E_1/2$, the distance is $\sqrt{E_1}=\sqrt{2E_b}$, and
$P_b=\Q\big(\sqrt{E_b/N_0}\big)$---the same as orthogonal, 3\,dB worse than antipodal. In terms
of \emph{peak} energy it is 6\,dB worse than antipodal, which matters for peak-power-limited
transmitters such as lasers.
\end{itemize}

\mdcfig[0.78]{ch08_ber_binary}{Bit error probability against $E_b/N_0$. Antipodal signalling
is 3\,dB better than orthogonal or on-off signalling at every error rate. The dots are a
Monte Carlo simulation of polar NRZ with $2\times10^6$ bits per point. Gray-coded 4-PAM needs
about 4\,dB more $E_b/N_0$ than binary; noncoherent detection of orthogonal FSK costs
roughly another 1\,dB at low error rates (Chapter~\ref{ch:modulation}).}

Figure~\ref{fig:ch08_ber_binary} plots these curves. The 3\,dB gap between antipodal and
orthogonal is a direct consequence of geometry: antipodal points are $2\sqrt{E_b}$ apart,
orthogonal points only $\sqrt2\sqrt{E_b}$. Everything else in digital modulation---QPSK, QAM,
coding gains---can be understood as a way of placing points so that their minimum distance is
as large as possible for a given average energy.

\begin{worked}[How much signal does a 10\,Gb/s receiver need?]
A 10\,Gb/s NRZ optical receiver must achieve a bit error rate of $10^{-12}$ without FEC. What
$E_b/N_0$ does ideal antipodal detection require, and what does that mean for the eye?

From~\eqref{eq:ch08:antipodal}, $\sqrt{2E_b/N_0}=\Q^{-1}(10^{-12})=7.03$, so
$E_b/N_0=7.03^2/2=24.7$, or 13.9\,dB. In terms of the eye: the ratio of half the eye opening to
the RMS noise, $d/2\sigma$, must be at least 7.03, or equivalently the full eye opening must be
about 14 times the RMS noise. If the receiver's input-referred noise after the matched filter is
$\sigma=2$\,mV, the eye must be at least 28\,mV high at the slicer---before accounting for any
ISI, jitter, crosstalk or slicer offset, each of which eats part of the margin. An optical
receiver's noise is actually signal dependent (shot noise and optical amplifier noise are larger
on a one than on a zero), which shifts the optimum threshold below the midpoint, a
practical instance of~\eqref{eq:ch08:map} with unequal variances rather than unequal priors.
Optical engineers express the same requirement as a ``Q-factor'' of 7.03, or 17\,dB in their
$20\log_{10}Q$ convention.
\end{worked}

\begin{inpractice}[The Q-factor and the BER that nobody measures]
A BER of $10^{-12}$ at 10\,Gb/s means one error every 100 seconds; measuring it with 95\%
confidence requires about $3\times10^{12}$ error-free bits, five minutes per data point, and
$10^{-15}$ would take days. Test engineers therefore extrapolate. They measure the mean and
standard deviation of the one and zero levels at the sampling instant (or the BER at several
deliberately offset thresholds, which is quicker to accumulate) and compute
$Q=(\mu_1-\mu_0)/(\sigma_1+\sigma_0)$, from which $\mathrm{BER}\approx\Q(Q)$. The same idea,
applied along the time axis, gives the bathtub curves of Section~\ref{sec:ch08:jitter}.
Extrapolation assumes Gaussian tails; bounded interferers such as crosstalk and deterministic
jitter make the true tails thinner, whereas rare events such as supply glitches make them
fatter, and the latter can surprise you in the field.
\end{inpractice}

%======================================================================
\section{Eye diagrams}
\label{sec:ch08:eyes}
%======================================================================

\subsection{Anatomy of an eye}

An \term{eye diagram} overlays many short segments of the received waveform, each two or three
symbols long, aligned to the symbol clock. On an oscilloscope it is produced by triggering on a
recovered clock and using infinite persistence; in simulation it is a reshaping of the waveform
array. Because every data pattern leaves its trace, the eye shows at a glance the worst-case
behaviour of the link: the effect of all combinations of neighbouring symbols, of noise and of
timing variations. It is the single most useful diagnostic in baseband engineering, and a
test engineer can read a link's health from it in seconds.

\mdcfig[0.9]{ch08_eye_anatomy}{Anatomy of an eye diagram (binary RRC signal, $\beta=0.5$,
$E_s/N_0=21$\,dB). The eye height is the vertical noise margin at the best sampling instant; the
eye width is the timing margin at the decision threshold. The slope of the sides determines how
quickly the vertical opening shrinks with timing error. A compliance mask (orange) defines a
keep-out region that no trace may enter.}

The features in Figure~\ref{fig:ch08_eye_anatomy} each tell a story:
\begin{itemize}
\item \emph{Eye height} (vertical opening) at the sampling instant is the noise margin: the
smallest distance between any ``one'' trace and any ``zero'' trace. ISI reduces it
deterministically (the spread of the upper and lower rails in the noise-free eye); noise
reduces it statistically.
\item \emph{Eye width} (horizontal opening) at the threshold is the timing margin, reduced by
jitter and by ISI, which shifts the zero crossings according to the data pattern.
\item \emph{Slope} of the eye's sides at the sampling instant sets the sensitivity to timing
error: a steep eye closes quickly as the sampling phase moves.
\item \emph{Crossing point}. For NRZ the crossings should be at 50\% of the amplitude. Crossings
that are high or low indicate asymmetric rise and fall times or duty-cycle distortion, typical
of optical transmitters.
\item \emph{Overshoot and ringing} indicate reflections from impedance discontinuities or an
over-peaked equaliser.
\end{itemize}

\subsection{Compliance masks}

Standards specify transmitter and receiver quality by \term{eye masks}: a hexagonal or
diamond-shaped keep-out region in the centre of the eye, and keep-out regions above and below,
that no trace may enter over a specified number of waveforms. SONET/SDH (ITU-T G.957) and
10G Ethernet optical transmitters are specified this way, with the mask margin (how far the
mask can be scaled before hits occur) as a figure of merit. Faster electrical standards have
moved from fixed masks to statistical ones: the receiver-side eye after a reference equaliser
must exceed a minimum eye height and width at a stated BER, obtained by extrapolation. PCI
Express, for example, defines the received eye after a reference CTLE and DFE at a BER of
$10^{-12}$; IEEE 802.3 PAM-4 interfaces use metrics such as \emph{eye height and width at
$10^{-6}$} (EH6, EW6) and the transmitter dispersion eye closure for PAM-4 (TDECQ), which
estimates the SNR penalty of a real transmitter compared with an ideal one after a reference
equaliser.

\subsection{Jitter and its decomposition}
\label{sec:ch08:jitter}

\textbf{Jitter}\index{jitter} is the deviation of signal edges from their ideal positions in time. At 25\,Gb/s
one unit interval (UI, a bit period) is 40\,ps, and a few picoseconds of jitter is a large
fraction of the budget. Jitter is decomposed into components with different statistics,
because they add differently:
\begin{itemize}
\item \textbf{Random jitter}\index{random jitter} (RJ): unbounded, approximately Gaussian, from thermal noise and
oscillator phase noise (Chapter~\ref{ch:sdr}). Specified by its RMS value $\sigma_{RJ}$.
\item \textbf{Deterministic jitter}\index{deterministic jitter} (DJ): bounded, specified by its peak-to-peak value. Its main
sub-types are \emph{data-dependent jitter} (DDJ), which is ISI seen on the time axis---an edge
after a long run of ones arrives at a different time from one after an alternating
pattern---\emph{duty-cycle distortion} (DCD), a systematic difference between rising and falling
edges; \emph{periodic jitter} (PJ) from power-supply ripple or spurious tones; and
\emph{bounded uncorrelated jitter} from crosstalk.
\end{itemize}
Bounded components add linearly (peak to peak), Gaussian components add in root-sum-square,
and the total jitter at a given BER is the combination:
\begin{equation}
\mathrm{TJ}(\mathrm{BER})\approx\mathrm{DJ}_{\delta\delta}+2\,\Q^{-1}\!\left(\frac{\mathrm{BER}}{\rho_T}\right)\sigma_{RJ},
\label{eq:ch08:tj}
\end{equation}
where $\rho_T$ is the transition density (0.5 for random data). This is the \term{dual-Dirac
model}: the DJ distribution is replaced by two impulses separated by $\mathrm{DJ}_{\delta\delta}$,
each convolved with the RJ Gaussian (Figure~\ref{fig:ch08_bathtub}, left). The factor multiplying
$\sigma_{RJ}$ is about 14 at $10^{-12}$ (13.9 with $\rho_T=0.5$; the often-quoted 14.07
assumes $\rho_T=1$) and about 16 at $10^{-15}$, so even a small RJ dominates
the total at low BER. $\mathrm{DJ}_{\delta\delta}$ is a model parameter fitted to the tails, and
it is usually smaller than the actual peak-to-peak DJ.

\subsection{Bathtub curves}

If the receiver samples at phase $x$ (in UI) from the left crossing of the eye, an error occurs
when an edge from the left crosses to the right of $x$, or one from the right crosses to its
left. Under the dual-Dirac model,
\begin{equation}
\mathrm{BER}(x)=\frac{\rho_T}{2}\left[\Q\!\left(\frac{x-\mathrm{DJ}/2}{\sigma}\right)+\Q\!\left(\frac{x+\mathrm{DJ}/2}{\sigma}\right)\right]+\big(\text{same with }x\to1-x\big).
\end{equation}
Plotted on a log scale against $x$, this is the \term{bathtub curve}
(Figure~\ref{fig:ch08_bathtub}, right): steep walls where the sampling phase is close to the
crossings, and a flat, very low floor in the middle. The width of the bathtub at a target BER is
the eye opening at that BER, $1-\mathrm{TJ}(\mathrm{BER})$. A bit-error-rate tester measures the
bathtub directly by sweeping the sampling phase and counting errors down to $10^{-8}$ or so;
the rest of the curve is extrapolated by fitting the dual-Dirac model to the measured walls.

\mdcfig{ch08_bathtub}{Left: a jitter histogram with a bimodal deterministic component convolved
with Gaussian random jitter; the RJ-only distribution (pink) is shown for comparison. Right:
bathtub curves for the dual-Dirac model. Increasing RJ from 0.01 to 0.035\,UI closes the
$10^{-12}$ eye far more than increasing DJ from 0.1 to 0.25\,UI, because RJ is multiplied by
about 14 at that error rate.}

\begin{worked}[A jitter budget at 25.78\,Gb/s]
A 25GBASE-R lane runs at 25.78125\,Gb/s, so 1\,UI $=38.8$\,ps. Suppose a transmitter has
$\mathrm{DJ}_{\delta\delta}=0.15$\,UI and $\sigma_{RJ}=0.7$\,ps $=0.018$\,UI. What is the total
jitter and eye width at $10^{-12}$?

From~\eqref{eq:ch08:tj} with $\rho_T=0.5$: $\Q^{-1}(2\times10^{-12})\approx6.93$, so
$\mathrm{TJ}=0.15+2(6.93)(0.018)=0.15+0.25=0.40$\,UI, and the eye width is $0.60$\,UI
$=23$\,ps. Notice that the 0.7\,ps of random jitter contributes more than the 5.8\,ps of
deterministic jitter. Halving $\sigma_{RJ}$ with a lower-phase-noise PLL would widen the eye by
0.125\,UI (4.8\,ps); halving the DJ would add only 0.075\,UI. The lesson generalises: at low
BER, random jitter is the expensive kind.
\end{worked}

\subsection{Clock and data recovery}

Every eye in this chapter assumed that the receiver knows where the centre of the eye is. In a
wireline link no clock is sent; the receiver derives it from the data with a
\term{clock and data recovery} (CDR) circuit, a phase-locked loop whose phase detector compares
data transitions with the local clock. The classic \emph{Alexander} or \emph{bang-bang} phase
detector samples each bit twice, at the centre and at the edge: if the edge sample agrees with
the previous bit, the clock is late; if it agrees with the next bit, early. Baud-rate detectors
such as the Mueller--M\"uller detector use only one sample per symbol and are favoured in
ADC-based receivers. The CDR's loop bandwidth is a compromise (many serial standards
measure jitter through a reference ``golden PLL'' with a corner at the baud rate divided by
1667, which is 3.75\,MHz at 6.25\,Gb/s): it must track
low-frequency wander and spread-spectrum clocking, which the receiver tolerates, while rejecting
high-frequency jitter, which it cannot follow and which therefore eats the eye. The
standards express this as a \emph{jitter tolerance mask}: sinusoidal jitter of many UI is
tolerated at low frequency, falling to a fraction of a UI above the loop bandwidth.
Chapter~\ref{ch:sync} develops the PLL theory and timing-error detectors in detail.

%======================================================================
\section{Channel loss and baseband equalisation}
\label{sec:ch08:channel}
%======================================================================

So far the channel has been flat. Real baseband channels---a circuit-board trace, a backplane,
a twinax cable, a twisted pair---are not. Their loss rises steeply with frequency, and at modern
data rates the received pulse is spread over many symbol periods. This section introduces the
problem and the three classic remedies; Chapter~\ref{ch:equalization} develops the theory in
full.

\subsection{Why copper is a low-pass filter}

Two mechanisms dominate the loss of a transmission line at gigahertz frequencies. The
\term{skin effect} confines current to a thin layer at the conductor surface, of depth
$\delta=1/\sqrt{\pi f\mu\sigma}$ (about 0.66\,$\mu$m in copper at 10\,GHz), so resistance, and
hence loss in decibels, grows as $\sqrt f$. \textbf{Dielectric loss}\index{dielectric loss}, from the energy absorbed by
the insulating material in each cycle of the field, grows in proportion to $f$ and to the
material's loss tangent: about 0.02 for ordinary FR-4 laminate, and below 0.004 for the
low-loss materials used at 50\,Gb/s per lane and above. Surface roughness of the copper foil
adds more loss. On top of the smooth loss, every via stub, connector and package transition
causes a reflection, which appears as notches in the frequency response and echoes in the
pulse response. The engineering answers include low-loss laminates, smooth copper, back-drilled
vias to remove stubs, and ever-shorter paths such as cables in place of board traces and
co-packaged optics.

\mdcfig{ch08_channel_loss}{Left: a lossy-line model with loss in decibels proportional to
$0.4\sqrt f+0.6f$, scaled to 10, 20 and 30\,dB at the Nyquist frequency. Right: the resulting
single-bit (pulse) responses, normalised to their peaks. The dots mark the samples at
unit-interval spacing: the main cursor $h_0$, a small pre-cursor $h_{-1}$, and a long train of
post-cursors that are the ISI.}

The \term{pulse response}---the response of the channel to a single one-UI pulse---is the
wireline engineer's main tool. Sampled once per UI, it gives the \term{cursors}: the main
cursor $h_0$, pre-cursors $h_{-1},h_{-2},\dots$ and post-cursors $h_1,h_2,\dots$
(Figure~\ref{fig:ch08_channel_loss}). The worst-case eye opening can be computed directly from
them: the eye closes when the sum of the magnitudes of all the other cursors exceeds the main
cursor, $\sum_{k\ne0}|h_k|\ge h_0$. This \emph{peak distortion} analysis is how link budgets are
first estimated. Channel loss is quoted at the Nyquist frequency, $f_N=R_s/2$: in rough terms, a
link with less than 10\,dB of loss at $f_N$ needs little equalisation, one with 20--30\,dB needs
serious equalisation, and the long-reach backplane channels of current standards---about
28\,dB at 26.56\,GHz for IEEE~802.3ck 100GBASE-KR1, about 36\,dB at 16\,GHz for PCI Express
5.0---need everything in the next subsection.

\subsection{Three equalisers}

\begin{figure}[htbp]\centering
\begin{tikzpicture}[node distance=0.35cm and 0.38cm,
  blk/.style={draw=navy,fill=navy!6,thick,minimum height=0.95cm,minimum width=1.2cm,align=center,font=\sffamily\scriptsize},
  arr/.style={-{Stealth[length=2mm]},thick,navy}]
\node[blk] (ser) {Serialiser\\(PISO)};
\node[blk,right=of ser] (ffe) {TX FFE\\(3--5 taps)};
\node[blk,right=of ffe] (drv) {Driver};
\node[blk,right=of drv,fill=accent!6,draw=accent,minimum width=1.6cm] (ch) {Channel\\package, board,\\connector, cable};
\node[blk,right=of ch] (ctle) {CTLE\\+ VGA};
\node[draw=accent,circle,thick,right=of ctle,minimum size=0.55cm,font=\small] (sum) {$+$};
\node[blk,right=of sum] (sl) {Slicer};
\node[blk,below=0.55cm of sl] (dfe) {DFE\\taps};
\node[blk,above=0.5cm of sl] (cdr) {CDR};
\draw[arr] (ser) -- (ffe); \draw[arr] (ffe) -- (drv); \draw[arr] (drv) -- (ch); \draw[arr] (ch) -- (ctle);
\draw[arr] (ctle) -- (sum); \draw[arr] (sum) -- (sl);
\draw[arr] (sl.east) -- ++(0.35,0) node[right,font=\sffamily\scriptsize]{bits};
\draw[arr] (sl.east) ++(0.15,0) |- (dfe.east);
\draw[arr] (dfe.west) -| node[pos=0.8,left,font=\scriptsize]{$-$} (sum.south);
\draw[arr,dashed] ([xshift=0.3cm]sl.north) -- node[right,font=\sffamily\scriptsize]{samples} ([xshift=0.3cm]cdr.south);
\draw[arr,dashed] ([xshift=-0.3cm]cdr.south) -- node[left,font=\sffamily\scriptsize]{clock} ([xshift=-0.3cm]sl.north);
\end{tikzpicture}
\caption{A typical SerDes (serialiser--deserialiser) link. The transmitter pre-distorts with a
short feed-forward equaliser, the receiver's continuous-time linear equaliser boosts high
frequencies, and a decision-feedback equaliser subtracts the post-cursor ISI of past decisions.
The clock and data recovery (CDR) loop locks the sampling phase to the data. Above about
50\,Gb/s per lane, the receiver is often an ADC followed by a digital FFE and DFE instead.}
\label{fig:ch08_serdes}
\end{figure}

Figure~\ref{fig:ch08_serdes} shows the architecture used by nearly every high-speed serial
link. Three kinds of equaliser share the work.
\begin{description}[style=nextline,leftmargin=1.2em,font=\sffamily\bfseries\small\color{navy}]
\item[Feed-forward equaliser (FFE).] A finite-impulse-response filter with taps spaced one UI
apart. At the transmitter it is usually three to five taps: a large main tap with a negative
post-cursor tap is called \term{de-emphasis}, because it reduces the amplitude of bits that
repeat the previous one, and a negative pre-cursor tap is \emph{pre-shoot}. PCI Express 3.0
and later define a set of transmitter presets and let the receiver request adjustments during
link training. A transmit FFE cannot increase the peak voltage, so its high-frequency boost is
really a low-frequency cut. At the receiver, in ADC-based designs, the FFE may have 20 to 40
taps. Figure~\ref{fig:ch08_equalized_eye} shows a 9-tap FFE opening an eye that a 20\,dB channel
had closed completely.
\item[Continuous-time linear equaliser (CTLE).] An analog filter with a zero and two or more
poles, giving a high-pass ``peaking'' response of up to 10 to 20\,dB at the Nyquist frequency.
It is cheap, fast and power efficient, but it amplifies high-frequency noise and crosstalk
along with the signal.
\item[Decision-feedback equaliser (DFE).] Once a bit has been decided, its contribution to later
samples---the post-cursors $h_1,h_2,\dots$ times the decision---is known and can simply be
subtracted. Because the DFE works on noiseless decisions, it removes ISI without amplifying
noise, which is its great virtue. Its weaknesses are that it cannot touch pre-cursors, that a
wrong decision subtracts the wrong value and makes further errors more likely (\term{error
propagation}), and that the first tap must be computed and fed back within one UI---10\,ps at
100\,Gb/s. Designers meet the last constraint with \emph{loop unrolling} (speculation): compute
both possible outcomes in advance and select one when the previous decision arrives.
\end{description}

\mdcfig{ch08_equalized_eye}{Equalisation at baseband. Left: NRZ after the 20\,dB channel of
Figure~\ref{fig:ch08_channel_loss}; the eye is closed. Centre: after a 9-tap, UI-spaced
zero-forcing FFE with two pre-cursor and six post-cursor taps, the eye is wide open (no noise
is added here; in practice the FFE's high-frequency gain amplifies noise and crosstalk too).
Right: the FFE taps, dominated by the main tap and a negative first post-cursor.}

The division of labour matters for noise: a linear equaliser that fully inverts a 30\,dB
loss also boosts the high-frequency noise and crosstalk by up to 30\,dB, whereas a DFE cancels
post-cursor ISI without that penalty. That is why nearly every long-reach SerDes uses a
modest CTLE and FFE to shape the pre-cursors and the first few post-cursors, followed by a DFE.
Chapter~\ref{ch:equalization} derives the zero-forcing, MMSE and DFE solutions, adaptive
algorithms such as sign--sign LMS (used in SerDes because it needs only comparators), and the
maximum-likelihood sequence detectors that some 200\,Gb/s-per-lane receivers now include.

\begin{inpractice}[Anatomy of a 112\,Gb/s SerDes]
A 112\,Gb/s PAM-4 SerDes (56\,GBd, 53.125\,GBd with the Ethernet rate plan) as used for
IEEE~802.3ck and OIF CEI-112G links typically has a transmitter with a 3- to 5-tap FFE built
from a DAC, and a receiver with a CTLE, a 6- to 8-bit ADC running at the symbol rate
(time-interleaved from 32 to 64 slower sub-ADCs), and a DSP with an FFE of a few tens of taps, a
one-tap or short DFE or a simple sequence detector, a CDR and adaptation engines. The whole
lane consumes a few hundred milliwatts, and a switch chip with 512 such lanes (51.2\,Tb/s)
spends a large part of its power budget on SerDes. This is why the industry is so interested in
shorter reaches, linear-drive optics and co-packaged optics, all of which reduce the loss a
SerDes must overcome.
\end{inpractice}

%======================================================================
\section{Multilevel PAM}
\label{sec:ch08:mpam}
%======================================================================

\subsection{\texorpdfstring{$M$}{M}-PAM and its error probability}

An $M$-level PAM alphabet places the levels at $\{\pm1,\pm3,\dots,\pm(M-1)\}\cdot d/2$, equally
spaced by $d$. Each symbol carries $\log_2M$ bits. With a matched filter and ML thresholds at the
midpoints, an inner level can be mistaken for either neighbour and an outer level for one, so
the symbol error probability is
\begin{equation}
P_s=2\left(1-\frac1M\right)\Q\!\left(\frac{d}{2\sigma}\right).
\end{equation}
The average energy is $E_s=\frac{(M^2-1)}{12}d^2$ (for unit-energy pulses), and
$E_b=E_s/\log_2M$. Substituting $\sigma^2=N_0/2$ gives
\begin{equation}
P_s=2\left(1-\frac1M\right)\Q\!\left(\sqrt{\frac{6\log_2M}{M^2-1}\cdot\frac{E_b}{N_0}}\right).
\label{eq:ch08:mpam}
\end{equation}
With \term{Gray coding}---labelling adjacent levels with bit patterns that differ in one bit, as
in the PAM-4 map $\{00,01,11,10\}$---a symbol error almost always causes exactly one bit
error, so $P_b\approx P_s/\log_2M$.

The factor $6\log_2M/(M^2-1)$ is 2 for binary, 0.8 for PAM-4, $12/35\approx0.34$ for PAM-8.
At a fixed $E_b/N_0$, PAM-4 therefore loses $10\log_{10}(2/0.8)\approx4$\,dB relative to
binary, and each further doubling of $M$ costs about 6\,dB more. In return, PAM-4 carries
twice the bit rate at the same symbol rate, or the same bit rate in half the bandwidth.
Chapter~\ref{ch:infotheory} shows that this exchange of power for bandwidth is exactly what
Shannon's capacity formula predicts in the high-SNR regime.

\mdcfig{ch08_pam4}{NRZ and PAM-4 at the same bit rate and the same peak swing. PAM-4 runs at
half the symbol rate, so its eyes are twice as wide in absolute time, but each of its three
eyes is one third as tall. The Gray-coded labels are shown at right. Transitions between the
outer levels ($00\leftrightarrow10$) cross all three thresholds and are the steepest, which is
why PAM-4 eyes look skewed in real transmitters with limited slew rate.}

\subsection{Why the data centre switched to PAM-4}

For a band-limited wireless channel with flat noise, the 4\,dB penalty is simply the price of
spectral efficiency. For a lossy copper channel it is a bargain, because halving the symbol
rate halves the Nyquist frequency, and the loss at the Nyquist frequency falls by far more than
4\,dB. There is a second way to state the penalty that is common in the SerDes world: at a fixed
\emph{peak-to-peak swing} (set by the supply voltage and the driver), each PAM-4 eye is one third
of the NRZ eye, a loss of $20\log_{10}3=9.5$\,dB in vertical opening
(Figure~\ref{fig:ch08_pam4}). Part of that is regained because the noise bandwidth is halved.

\begin{worked}[NRZ or PAM-4 at 112\,Gb/s?]
A backplane channel has a loss of approximately 0.9\,dB/GHz in the 20--60\,GHz range. Compare NRZ
and PAM-4 at 112\,Gb/s.

NRZ runs at 112\,GBd with a Nyquist frequency of 56\,GHz, where the loss is about
$0.9\times56\approx50$\,dB. PAM-4 runs at 56\,GBd with a Nyquist frequency of 28\,GHz and a loss of
about 25\,dB. The PAM-4 signal starts with a 9.5\,dB eye-height disadvantage, but arrives with
25\,dB less loss to equalise; since equaliser noise enhancement and residual ISI grow rapidly with
the loss being compensated, PAM-4 wins by a large margin. A 50\,dB NRZ channel is, for practical
purposes, not equalisable with SerDes-class power budgets. Conversely, on a short, low-loss link
(a few decibels of loss), NRZ is better, which is why NRZ survived at 25\,Gb/s and PAM-4 took over
at 50 and 100\,Gb/s per lane.
\end{worked}

\subsection{PAM-4 in standards}

PAM-4 has become the workhorse of the data centre. Table~\ref{tab:ch08:pam4} lists the main
deployments.

\begin{table}[htbp]\centering\small
\caption{Multilevel PAM in wireline standards (rates are per lane; approximate years of
publication).}
\label{tab:ch08:pam4}
\begin{tabularx}{\linewidth}{@{}l l l X@{}}
\toprule
Standard & Year & Signalling & Notes\\
\midrule
IEEE 802.3bs (200/400GbE) & 2017 & 26.5625\,GBd PAM-4 electrical & 400GAUI-8; RS(544,514) ``KP4'' FEC mandatory\\
IEEE 802.3cd (50/100/200GbE) & 2018 & 26.5625\,GBd PAM-4 & 50\,Gb/s per lane over backplane and copper cable\\
IEEE 802.3ck & 2022 & 53.125\,GBd PAM-4 & 100\,Gb/s per electrical lane; basis of 400G and 800G ports\\
IEEE 802.3dj & c.\ 2026 & 106.25\,GBd PAM-4 & 200\,Gb/s per lane, for 800G and 1.6T Ethernet; optical lanes with an inner FEC run at about 113\,GBd\\
PCI Express 6.0 & 2022 & 32\,GBd PAM-4 (64\,GT/s) & Flit mode, light-weight FEC plus CRC, no block code\\
PCI Express 7.0 & 2025 & 64\,GBd PAM-4 (128\,GT/s) & Same flit architecture\\
GDDR6X (Micron) & 2020 & PAM-4, about 10\,GBd per pin & 19--21\,Gb/s per pin in graphics cards\\
GDDR7 (JEDEC) & 2024 & PAM-3 & about 1.5 bits per symbol, simpler than PAM-4\\
USB4 Version 2.0 & 2022 & PAM-3 & 80\,Gb/s, 11 bits in 7 ternary symbols\\
\bottomrule
\end{tabularx}
\end{table}

Three design features recur in PAM-4 systems.
\begin{itemize}
\item \emph{Forward error correction is mandatory.} Without FEC, PAM-4 links could not reach a
BER of $10^{-12}$ at acceptable power. Ethernet's KP4 code, the Reed--Solomon RS(544,514) code
over 10-bit symbols (Chapter~\ref{ch:classiccodes}), corrects up to 15 symbol errors per
codeword and turns a raw BER of about $2.4\times10^{-4}$ into an essentially error-free link, at
the cost of about 100\,ns of latency. PCI Express 6.0, which cannot afford that latency, uses a
light-weight interleaved code and relies on a CRC and fast link-level retry, with a target raw
flit error rate of around $10^{-6}$.
\item \emph{Precoding against error bursts.} A DFE error on a PAM-4 signal tends to cause a
burst of errors. The optional $1/(1+D)$ modulo-4 precoder in the Ethernet standards converts
such bursts into (typically) two symbol errors, which the Reed--Solomon code handles easily.
This is the same precoding idea as in duobinary (Section~\ref{sec:ch08:pr}).
\item \emph{Level linearity.} The three eyes are equal only if the four levels are equally
spaced. Transmitters are specified by a \emph{ratio of level mismatch} (RLM), typically at least
0.95, and optical PAM-4 transmitters, whose lasers and modulators are nonlinear, are often
pre-distorted to space the levels. GDDR6X adds a coding trick, \emph{maximum transition
avoidance}, that forbids full-swing transitions between the outer levels to reduce
simultaneous-switching noise on wide parallel buses.
\end{itemize}

\begin{inpractice}[PAM-4 in fibre]
Direct-detection optical links also use PAM-4: 400GBASE-DR4 sends four lanes of 53.125\,GBd
PAM-4 (100\,Gb/s each) over single-mode fibre up to 500\,m, and 800G modules use eight such
lanes or four lanes at twice the rate. An optical PAM-4 receiver sees intensity, not field, so
the levels are optical power levels and the noise differs between them; the TDECQ metric of
Section~\ref{sec:ch08:eyes} was invented for exactly this. For longer reaches, where chromatic
dispersion and the square-law detector make intensity modulation impractical, coherent optics
with QAM on both polarisations takes over (Chapter~\ref{ch:wireline}).
\end{inpractice}

%======================================================================
\section{Partial-response signalling}
\label{sec:ch08:pr}
%======================================================================

The Nyquist criterion treats ISI as the enemy. \textbf{Partial-response signalling}\index{partial-response signalling} (also called
correlative coding) treats it as a design variable: introduce a small, \emph{known} amount of ISI
on purpose, and in exchange obtain a pulse that is realisable at exactly the Nyquist bandwidth,
or a spectrum with nulls where the channel has none to give.

\subsection{Duobinary}

The \term{duobinary} pulse is the sum of two sinc pulses one symbol apart:
\begin{equation}
x(t)=\sinc\!\left(\frac tT\right)+\sinc\!\left(\frac{t-T}{T}\right),\qquad
X(f)=\begin{cases}2T\cos(\pi fT)\,e^{-j\pi fT},&|f|<\dfrac1{2T},\\0,&\text{otherwise.}\end{cases}
\label{eq:ch08:duobinary}
\end{equation}
Its spectrum occupies exactly the Nyquist bandwidth $1/(2T)$---no excess bandwidth at all---yet,
unlike the sinc, it falls smoothly to zero at the band edge, so it can be approximated by a
realisable filter (Figure~\ref{fig:ch08_duobinary}). Its tails decay as $1/t^2$. The price is
that $x(t)$ is 1 at \emph{two} sampling instants, $t=0$ and $t=T$. The sampled output is
\begin{equation}
y_k=a_k+a_{k-1},
\end{equation}
which, for binary $a_k=\pm1$, takes three values: $+2$, $0$ and $-2$. The ISI is controlled: the
receiver knows it is exactly the previous symbol. In $D$-transform notation (with $D$ a one-symbol
delay), the channel from symbols to samples is $1+D$.

\mdcfig{ch08_duobinary}{Partial response. Left: the duobinary spectrum ($1+D$) falls smoothly
to zero at the Nyquist frequency; modified duobinary ($1-D^2$) also has a null at DC. Centre:
the duobinary pulse is the sum of two sincs and is non-zero at exactly two sampling instants.
Right: the sampled duobinary signal has three levels and two eyes.}

A receiver could decide $a_k$ as $\hat a_k=y_k-\hat a_{k-1}$, but one error would then propagate
indefinitely. The standard cure is \term{precoding}, applied at the transmitter to the bits
$d_k\in\{0,1\}$:
\begin{equation}
b_k=d_k\oplus b_{k-1},\qquad a_k=2b_k-1.
\end{equation}
Then $y_k=a_k+a_{k-1}=2(b_k+b_{k-1}-1)$ is 0 exactly when $b_k\ne b_{k-1}$, that is, when
$d_k=1$, and $\pm2$ when $d_k=0$. The receiver decides each bit from a single sample---$\hat d_k=1$
if $|y_k|<1$, else 0---with no memory and no error propagation. Precoding moves the inverse of
the channel's memory to the transmitter, where there is no noise; the same principle reappears
in Tomlinson--Harashima precoding and in the precoders of multi-user MIMO
(Chapter~\ref{ch:mimo}).

\begin{worked}[Precoded duobinary by hand]
Encode the bits $d=1\,0\,1\,1\,0\,0\,1$ with precoded duobinary, starting from $b_{-1}=0$.

The precoder gives $b_k=d_k\oplus b_{k-1}$: $b=1\,1\,0\,1\,1\,1\,0$. The levels are
$a=+1,+1,-1,+1,+1,+1,-1$, and with $a_{-1}=-1$ the channel outputs $y_k=a_k+a_{k-1}$ are
$0,+2,0,0,+2,+2,0$. Decoding with ``$|y|<1\Rightarrow1$'' gives $1\,0\,1\,1\,0\,0\,1$, the
original data. Now suppose noise corrupts $y_2$ from 0 to $+2$. Only $\hat d_2$ is wrong; every
other bit is decoded independently. Without precoding, the recursive decision
$\hat a_k=y_k-\hat a_{k-1}$ would carry the error forward.
\end{worked}

\subsection{The cost, and how to recover it}

Symbol-by-symbol detection of duobinary costs noise margin. The receiver now makes a
three-level decision, and each of the two eyes spans only half the peak-to-peak range of the
signal, whereas a binary eye spans all of it. When the duobinary response is split optimally
between the transmit and receive filters, the loss works out to a factor of $(4/\pi)^2$, or
about 2.1\,dB of $E_b/N_0$, relative to ideal binary signalling with Nyquist pulses (see
Proakis and Salehi in the further reading). It can be recovered almost entirely by recognising that the
channel $1+D$ is a two-state machine and deciding the whole sequence at once with the
\term{Viterbi algorithm} (maximum-likelihood sequence detection, Chapters~\ref{ch:classiccodes}
and~\ref{ch:equalization}). The minimum distance between two different output sequences of the
$1+D$ channel equals that of binary signalling with the same pulse energy, so sequence detection
suffers no asymptotic loss.

\subsection{Modified duobinary and the partial-response classes}

Other polynomials give other spectra. In 1966 Kretzmer tabulated the useful ones as
\emph{classes} of partial response. The most important besides duobinary (class~1, $1+D$) is
\term{modified duobinary} (class~4, $1-D^2$), whose spectrum
$\propto\sin(2\pi fT)$ has nulls at \emph{both} DC and the Nyquist frequency
(Figure~\ref{fig:ch08_duobinary}). A DC null suits transformer-coupled lines and, as the next
subsection shows, magnetic recording. The precoder is $b_k=d_k\oplus b_{k-2}$, and the output
again has three levels. The $1-D^2$ channel is two interleaved $1-D$ channels (even and odd
samples), so its Viterbi detector is two interleaved two-state detectors, which made it cheap to
build in 1990s silicon.

\begin{history}[Duobinary, from telephone lines to 40\,Gb/s optics]
Adam Lender of Lenkurt Electric published the duobinary technique in 1963, motivated by the
need to send more data over telephone-grade analog channels. Ernest Kretzmer of Bell Labs
generalised it to partial-response classes in 1966. The idea found its most enduring homes far
from telephony. Optical duobinary, in which the three levels $-2,0,+2$ are sent as optical
fields of opposite phase so that the intensity has only two levels, was used in 10 and 40\,Gb/s
fibre systems from around 2000 because its narrow spectrum tolerates chromatic dispersion. And
class-4 partial response became, as ``PR4'', the target of the read channels in virtually every
hard disk drive of the 1990s.
\end{history}

\subsection{PRML: partial response in the disk drive}

A magnetic recording head reads transitions of magnetisation. Each transition produces a
pulse---well modelled by the Lorentzian $s(t)=1/[1+(2t/PW_{50})^2]$, where $PW_{50}$ is the pulse
width at half its amplitude---and a written bit, which is a pair of opposite transitions,
produces the \emph{dibit} response $s(t)-s(t-T)$. The channel therefore has a built-in $1-D$
factor, a null at DC (a head cannot read a constant magnetisation), and a high-frequency
roll-off that worsens as the \emph{normalised density} $PW_{50}/T$ increases
(Figure~\ref{fig:ch08_prml}). Early drives used peak detection: find each pulse's peak and
declare a transition. That works only if pulses are well separated, so densities were limited
to $PW_{50}/T$ below about 1, and run-length-limited codes such as (2,7) were used to keep
transitions apart.

\mdcfig{ch08_prml}{Left: the frequency response of a Lorentzian recording channel (dibit
response) at three normalised densities. Right: the partial-response targets used in
read channels. PR4 suits densities near 1.5--2; the extended targets EPR4 and E$^2$PR4, with more
low-frequency emphasis, match the higher densities better.}

\term{PRML} (partial response, maximum likelihood) takes the opposite approach. An equaliser
shapes the channel to a partial-response target such as PR4, $1-D^2=(1-D)(1+D)$, whose
spectrum matches the recording channel's natural shape (a DC null and a high-frequency
roll-off) so that little noise enhancement is needed; a Viterbi detector then decides the
sequence. The idea was proposed by Kobayashi and Tang of IBM in 1970, and IBM shipped the first
PRML disk drive around 1990. As densities rose, the targets were extended to EPR4,
$(1+D)(1-D^2)$, and E$^2$PR4, $(1+D)^2(1-D^2)$, whose Viterbi trellises have 8 and 16 states.
Later drives replaced fixed targets with noise-predictive maximum-likelihood (NPML) detectors
and, from the late 2000s, iterative detection with LDPC codes (Chapter~\ref{ch:moderncodes}). The
PRML read channel roughly doubled the density achievable with peak detection and was one of
the key enablers of the capacity explosion of hard drives in the 1990s and 2000s.

\begin{inpractice}[Partial response is everywhere once you look]
The duobinary idea keeps reappearing. The channel of Section~\ref{sec:ch08:channel}, with its
large first post-cursor, is naturally close to $1+D$; some 100 and 200\,Gb/s-per-lane
receivers equalise to a $1+\alpha D$ target and use a short Viterbi detector (``1+D MLSE'') rather
than forcing the full response to a Nyquist pulse, avoiding the noise enhancement of complete
equalisation. Tamed FM and GMSK in GSM (Chapter~\ref{ch:modulation}) can be viewed as
partial-response continuous-phase modulation. And the $1/(1+D)$ precoder in 802.3 PAM-4 links
is Lender's precoder generalised to four levels.
\end{inpractice}

%======================================================================
\section{Faster-than-Nyquist signalling}
\label{sec:ch08:ftn}
%======================================================================

The Nyquist rate of $2W$ symbols per second is the maximum rate for \emph{zero ISI}. It is not a
limit on the information rate, and it is not even a limit on the symbol rate. In
\term{faster-than-Nyquist} (FTN) signalling, pulses designed for a symbol period $T$ are sent
every $\tau T$ with $\tau<1$, without changing the pulse or the bandwidth:
\begin{equation}
s(t)=\sum_k a_k\,p(t-k\tau T),\qquad 0<\tau\le1.
\end{equation}
The pulses are no longer orthogonal and ISI is unavoidable (Figure~\ref{fig:ch08_ftn}, left), so
symbol-by-symbol detection fails. But the symbol rate is $1/\tau$ times higher in the same
bandwidth, and a sequence detector can undo the ISI.

The surprising result, found by James Mazo at Bell Labs in 1975, concerns the minimum Euclidean
distance between two different signals. For binary sinc pulses, the minimum distance stays at
the binary antipodal value $2\sqrt{E_b}$---so an ideal sequence detector has \emph{the same
asymptotic error rate as ISI-free binary signalling}---for all $\tau$ down to about 0.802. This
is the \term{Mazo limit}. Below it, some pair of long error events comes closer together and the
distance falls (Figure~\ref{fig:ch08_ftn}, right). Signalling at $\tau=0.802$ carries 25\% more
bits in the same bandwidth at no cost in asymptotic power efficiency. For raised-cosine pulses
with excess bandwidth, the Mazo limit is lower still, since there is more ``slack'' spectrum to
exploit.

\mdcfig{ch08_ftn}{Faster-than-Nyquist signalling. Left: sinc pulses sent every $0.8T$ no longer
have zero crossings at the other sampling instants. Right: the minimum squared distance between
binary FTN signals, normalised to the antipodal value $4E_b$, found by a search over error
events up to 8 bits long. It remains at 0\,dB down to $\tau\approx0.80$, the Mazo limit, and
falls below it.}

FTN does not beat Shannon: a Nyquist-rate system with a larger alphabet and good coding can
achieve the same rates. Its appeal is practical---binary or QPSK transmitters and power amplifiers
with the spectral efficiency of denser constellations---and its cost is receiver complexity, since
the ISI spans many symbols and full Viterbi detection is out of the question. Reduced-complexity
detectors (M-algorithm, BCJR with channel shortening, iterative detection combined with the
decoder) have made FTN practical in experiments and it has been studied for satellite links,
optical superchannels and 6G, a story taken up in Chapter~\ref{ch:sixg}. The deeper point is conceptual: the
Nyquist rate is a property of a particular receiver, not of the channel.

%======================================================================
\section{Summary}
%======================================================================

\begin{itemize}
\item A baseband PAM signal $\sum_ka_kp(t-kT)$ splits the design into the alphabet (bits per
symbol, noise tolerance) and the pulse (bandwidth, ISI, timing sensitivity). Bandwidth follows
the symbol rate, measured in baud.
\item Line codes adapt the signal to real channels: DC balance for AC coupling, transitions for
clock recovery, spectral shaping, error monitoring and framing. The classic bit-level codes
(Manchester, AMI with B8ZS/HDB3, MLT-3, 2B1Q) have given way to block codes (8b/10b) and then to
scrambled, low-overhead codes (64b/66b, 128b/130b) and, with PAM-4 and FEC, to no line code at
all.
\item The PSD of a PAM signal is $|P(f)|^2S_a(f)/T$. A non-zero symbol mean adds spectral lines;
correlation between symbols shapes the spectrum and can put nulls at DC.
\item Zero ISI requires the folded spectrum $\sum_kX(f-k/T)$ to be flat. The minimum bandwidth is
$R_s/2$; the raised-cosine family trades excess bandwidth $\beta$ for shorter filters, timing
tolerance and lower peak-to-average ratio. Split as a root-raised-cosine pair, it is both
Nyquist and matched.
\item The matched filter, $h(t)=p^*(T_0-t)$, maximises SNR at $2E_p/N_0$, independent of pulse
shape; with coloured noise, whiten first. Antipodal signalling gives $P_b=\Q(\sqrt{2E_b/N_0})$,
3\,dB better than orthogonal or on-off signalling; $M$-PAM trades about 6\,dB per extra bit per
symbol for bandwidth.
\item Eye diagrams, jitter decomposition and bathtub curves are how links are measured. Random
jitter is multiplied by about 14 at $10^{-12}$, so it dominates budgets at low BER.
\item Copper loss grows with frequency; FFE, CTLE and DFE equalisers remove ISI, with DFE avoiding
noise enhancement. Halving the Nyquist frequency is why PAM-4, with mandatory FEC, now carries
Ethernet, PCIe~6 and graphics memory.
\item Partial response introduces known ISI to reach the Nyquist bandwidth with realisable filters
or to match a channel's spectral shape; precoding removes error propagation and Viterbi detection
recovers the distance. PRML made the modern hard disk. Faster-than-Nyquist signalling shows that
even the Nyquist rate is a receiver convention, not a law.
\end{itemize}

\begin{labbox}[Pulse shaping, line codes and eye diagrams]
Two notebooks accompany this chapter:
\begin{itemize}
\item \lab{moderncomms-labs/labs/lab03\_pulse\_shaping.py}
\item \lab{moderncomms-labs/labs/lab13\_line\_codes\_eyes.py}
\end{itemize}
The first compares rectangular, sinc and raised-cosine pulses in time and frequency, verifies the Nyquist
criterion, demonstrates that the matched filter maximises SNR and why the RRC is split between
transmitter and receiver, and lets you read eye diagrams as you vary roll-off, noise and timing
offset (Figures~\ref{fig:ch08_raised_cosine}, \ref{fig:ch08_eyes}
and~\ref{fig:ch08_timing_sensitivity}). The second implements the line codes of Section~\ref{sec:ch08:linecodes} and measures their spectra against
the formulas of Section~\ref{sec:ch08:psd}, builds LFSR scramblers and an 8b/10b encoder and
plots the running disparity, draws eye diagrams and bathtub curves with adjustable random and
deterministic jitter, compares NRZ with PAM-4 over a lossy channel with a simple FFE, and encodes
and decodes precoded duobinary. The figure script \lab{book/figscripts/ch08\_figs.py} generates
every data figure in this chapter and is a useful starting point for your own experiments.
\end{labbox}

\problems
\begin{enumerate}[label=\thechapter.\arabic*]
\item A modem sends 3200 symbols per second using a 64-point constellation. What are its bit
rate and baud rate? What minimum bandwidth would a baseband PAM signal need at the same symbol
rate?
\item Sketch the NRZ-L, NRZI, Manchester, AMI and MLT-3 waveforms for the bit sequence
\texttt{1100001011}. Which of them have a transition in every bit? Which are DC-free for this
sequence?
\item Encode the T1 sequence \texttt{1 0000 0000 1 0000 0000 01} with B8ZS, assuming the last
mark before it was positive. Mark the bipolar violations. Repeat with HDB3 for E1.
\item Using~\eqref{eq:ch08:psdlines}, derive the PSD of unipolar RZ signalling with a pulse of
duty cycle $\delta$ (pulse width $\delta T$). What fraction of the total power is in the
spectral line at $f=1/T$? Which $\delta$ maximises it?
\item Show that AMI with independent equiprobable input bits has $R_a[0]=\tfrac12$,
$R_a[\pm1]=-\tfrac14$ and $R_a[m]=0$ for $|m|\ge2$, and hence $S_a(f)=\sin^2(\pi fT)$.
\item Prove that the raised-cosine spectrum~\eqref{eq:ch08:rcspec} satisfies the Nyquist
criterion~\eqref{eq:ch08:nyquist}. Then show that the trapezoidal spectrum (linear roll-off
over the same band) also does. Compare the decay rates of the two impulse responses.
\item A UMTS carrier uses RRC filtering with $\beta=0.22$ at 3.84\,Mchip/s. What is the
occupied bandwidth? If the symbol rate were raised until the occupied bandwidth filled the full
5\,MHz channel, what would it be?
\item Derive the matched filter for coloured noise,~\eqref{eq:ch08:coloured}, using the
Cauchy--Schwarz inequality in the frequency domain. Show that it can be implemented as a whitening
filter followed by a filter matched to the whitened pulse.
\item A binary source emits ones with probability 0.9. The receiver sees $y=\pm1+\nu$ with
$\sigma=0.5$. Find the MAP threshold and the error probability, and compare with the ML
threshold. At what $\sigma$ does the MAP threshold move by more than 0.1?
\item A 28\,Gb/s NRZ link has $\mathrm{DJ}_{\delta\delta}=0.2$\,UI. What RMS random jitter can it
tolerate if the eye width at $10^{-12}$ must be at least 0.3\,UI? Express the answer in
picoseconds.
\item Using~\eqref{eq:ch08:mpam}, find the $E_b/N_0$ required for a bit error rate of $10^{-6}$
with 2-, 4- and 8-PAM (Gray coded). How much of the penalty is regained if the channel has 1\,dB
of loss per GHz and the bit rate is 50\,Gb/s?
\item The sampled pulse response of a channel is $h=[0.05,\,1,\,0.45,\,0.2,\,0.1]$ (pre-cursor,
main, post-cursors). Compute the worst-case eye opening with no equaliser. Design (a) a three-tap
transmit FFE and (b) a three-tap DFE that remove the first three post-cursors, and compare the
remaining worst-case ISI. Which one amplifies noise?
\item Show that the $1-D^2$ modified-duobinary channel with precoder $b_k=d_k\oplus b_{k-2}$ can be
decoded symbol by symbol as $\hat d_k=1$ if $|y_k|>1$. Encode and decode a sequence of ten bits by
hand.
\item \emph{(Simulation)} Implement a 64b/66b scrambler and descrambler. Introduce a single bit
error on the line and verify that it produces exactly three errors after descrambling. Then show
that the descrambler resynchronises from an arbitrary initial state after 58 bits.
\item \emph{(Simulation)} Generate 10\,000 symbols of RRC-shaped binary PAM with $\beta=0.25$ and
a span of 8 symbols. Measure the residual ISI (noise-free EVM) and the adjacent-channel leakage
ratio. Repeat with spans of 16 and 32 and with a Kaiser window applied to the taps.
\item \emph{(Simulation)} Reproduce Figure~\ref{fig:ch08_equalized_eye} and extend it: add white
noise at the receiver, and compare the BER of (a) a 9-tap ZF FFE, (b) a 9-tap MMSE FFE and (c) a
3-tap FFE followed by a 5-tap DFE, as the channel loss increases from 10 to 30\,dB.
\end{enumerate}

\furtherreading
\begin{itemize}
\item J.~G.~Proakis and M.~Salehi, \emph{Digital Communications}, 5th ed., McGraw-Hill, 2008.
Chapter 9 covers band-limited channels, the Nyquist criterion, partial response and the
spectral analysis of PAM; the standard graduate reference.
\item E.~A.~Lee and D.~G.~Messerschmitt, \emph{Digital Communication}, 2nd ed., Kluwer, 1994. The
best treatment of baseband PAM, line codes, the whitened matched filter and equalisation from the
wireline point of view.
\item H.~Nyquist, ``Certain topics in telegraph transmission theory,'' \emph{Transactions of the
AIEE}, vol.~47, pp.~617--644, 1928 (reprinted in \emph{Proc. IEEE}, vol.~90, no.~2, 2002). Still
worth reading for its clarity.
\item A.~X.~Widmer and P.~A.~Franaszek, ``A DC-balanced, partitioned-block, 8B/10B transmission
code,'' \emph{IBM Journal of Research and Development}, vol.~27, no.~5, 1983. The original 8b/10b
paper, with the complete tables.
\item J.~Bellamy, \emph{Digital Telephony}, 3rd ed., Wiley, 2000. T1/E1 line coding, B8ZS, HDB3 and
the practice of the telephone plant.
\item J.~E.~Mazo, ``Faster-than-Nyquist signaling,'' \emph{Bell System Technical Journal},
vol.~54, no.~8, 1975; and J.~B.~Anderson, F.~Rusek and V.~\"Owall, ``Faster-than-Nyquist
signaling,'' \emph{Proceedings of the IEEE}, vol.~101, no.~8, 2013, a modern survey.
\item P.~Kabal and S.~Pasupathy, ``Partial-response signaling,'' \emph{IEEE Transactions on
Communications}, vol.~23, no.~9, 1975. The classic tutorial on duobinary, modified duobinary and
the partial-response classes.
\item R.~D.~Cideciyan, F.~Dolivo, R.~Hermann, W.~Hirt and W.~Schott, ``A PRML system for digital
magnetic recording,'' \emph{IEEE Journal on Selected Areas in Communications}, vol.~10, no.~1,
1992. How PRML was engineered into IBM's disk drives.
\item IEEE Std 802.3 (current edition), especially Clause 49 (64b/66b PCS), Clause 91 (RS-FEC) and
the PAM-4 clauses of 802.3bs, 802.3cd and 802.3ck. Dense, but the definitive source for SerDes
specifications, scramblers and compliance tests.
\item M.~P.~Li, \emph{Jitter, Noise, and Signal Integrity at High-Speed}, Prentice Hall, 2007. The
standard reference on jitter decomposition, the dual-Dirac model and bathtub curves.
\end{itemize}
